\documentclass[11pt]{article}
\usepackage[T1]{fontenc}
\usepackage[utf8]{inputenc}
\usepackage{lmodern,amsmath,amssymb,amsthm,mathtools,booktabs,longtable,array}
\usepackage[margin=1in]{geometry}
\usepackage[hyphens]{url}
\usepackage[colorlinks=true,linkcolor=blue,citecolor=blue,urlcolor=blue]{hyperref}
\usepackage{microtype}
\hypersetup{pdftitle={Decomposition-aware global optimization: certified coordinate grids, conditional recourse, and structural limits},pdfauthor={}}
\newtheorem{theorem}{Theorem}[section]
\newtheorem{lemma}[theorem]{Lemma}
\newtheorem{proposition}[theorem]{Proposition}
\newtheorem{corollary}[theorem]{Corollary}
\theoremstyle{remark}
\newtheorem{remark}[theorem]{Remark}
\newcommand{\R}{\mathbb R}
\newcommand{\Z}{\mathbb Z}
\newcommand{\E}{\mathbb E}
\newcommand{\norm}[1]{\lVert#1\rVert}
\newcommand{\dist}{\operatorname{dist}}
\newcommand{\poly}{\operatorname{poly}}
\newcommand{\OPT}{f^*}
\newcommand{\pathref}[2]{\href{run:#1}{\texttt{\detokenize{#2}}}}
\title{Decomposition-aware global optimization:\\
certified coordinate grids, conditional recourse, and structural limits}
\author{}
\date{}
\begin{document}
\maketitle
\begin{abstract}
Sparse factorization can make global optimization useful without making
every sparse problem easy. This report develops a complete certified grid
algorithm on supplied tree decompositions. A unary interpolation correction
makes finite-state dynamic programming a valid continuous lower bound;
coordinate min-marginals then remove regions that cannot improve a feasible
incumbent. Under global quadratic growth, the surviving geometric grids
have state counts independent of the requested accuracy. For rational
mixed-integer box quadratics this gives an exact algorithm with running
time $f(p,\kappa)\poly(I)$, where $p$ is maximum bag size, $I$ is binary
input length, and $\kappa$ is upper coordinate curvature divided by growth.
The growth constant is not required as input, and certificate validity does
not depend on it. Stationary-face recovery also gives finite exact output
for arbitrary nonunique rational box quadratics, without a favorable
general complexity bound. The approximation result extends to explicitly encoded
fixed-degree polynomials, with different exact-output obligations. We
separate these proved guarantees from implementation evidence and from
extensions that require additional recourse, geometric, or certification
structure. The report is a technical synthesis of repository results;
priority and competitive general-purpose performance are not asserted.
\end{abstract}
\tableofcontents

\section{The solver capability and its limits}\label{sec:scope}

The main capability is a global bound for a factorized objective on a
product domain, computed by finite tables that agree exactly on shared
coordinates. It can be used as a complete solver for that domain or as a
component of a larger constrained solver. The relevant parameter is
\emph{maximum bag size and conditioning together}. Small graph width
alone is not the claim. Large tables, weak growth, or a poor decomposition
can make the algorithm impractical even when the theorem applies.

The dependency chain starts with coordinate curvature, which yields a
corrected-grid lower bound and safe min-marginal filtering. Global growth
then gives localization, bounded state counts, and a quantitative output
guarantee.
The first two implications hold without uniqueness or growth. Growth is
used to bound work, not to justify the answer. Rational arithmetic and
quadratic height bounds convert the final implication into an exact
mixed-integer quadratic algorithm. Polynomial objectives require their
own output argument.

This distinction also places the result relative to earlier repository
work. The September 29 affine-message existence bounds were centered at
an optimizer and left an algorithmic issue on branching decompositions.
Shared-coordinate grids subsequently supplied an optimizer-independent
algorithm, and min-marginal filtering removed the accuracy exponent from
its state bound. This does not prove that an older overlapping-cell
affine-message representation has the same properties. In particular,
counterexamples based on separator-copy drift or accumulated interpolation
error still apply to the representations they concern.

The current external comparison also sets a sharp boundary. Del Pia and
Khajavirad give an exact $O(n^2)$ algorithm on forests, but prove strong
NP-hardness of continuous box QP at treewidth two, even with bounded
coefficients; quartic path problems are also strongly NP-hard
\cite{DelPiaKhajavirad2026}. The forest result needs no growth promise.
The conditioned theorem here therefore complements that work and does
not dominate it. Bienstock and Mu\~noz treat broader constrained polynomial
models with scaled feasibility and objective tolerances
\cite{BienstockMunoz2018}; the constrained extension below obtains exact
feasibility under narrower integral structure.

\section{Model, decomposition, and arithmetic}\label{sec:model}

Let
\begin{equation}\label{eq:model}
\min_{x\in X}F(x),\qquad F(x)=\sum_{a\in\mathcal A}f_a(x_{S_a}),
\qquad X=\prod_{i=1}^n X_i,
\end{equation}
where $X_i=[a_i,b_i]$ or $[a_i,b_i]\cap\Z$. Integer endpoints are
rounded inward, empty domains rejected, and fixed coordinates substituted
before the construction. The zero-dimensional case is immediate.
Thereafter $n\ge1$ and $s=\max_i(b_i-a_i)>0$.

A supplied tree $T$ has bags $B_t$, maximum size $p$, and $N$ nodes.
Every factor scope is contained in a bag to which the factor is assigned
once, and the bags containing any coordinate form a connected subtree.
The decomposition and assignments can be checked directly. The input
length $I$ includes these data. A variable can occur in arbitrarily many
bags; no occurrence bound is assumed.

A valid upper coordinate-curvature bound is $L>0$ such that
\begin{equation}\label{eq:semiconcavity}
t\longmapsto F(x_1,\ldots,t,\ldots,x_n)-\frac L2t^2
\quad\hbox{is concave on each coordinate interval.}
\end{equation}
The condition is on the full continuous hull, including fractional points
between integer labels, and on the \emph{sum} of all factors. For a quadratic
with Hessian $H$, $L=\max_i\max\{H_{ii},0\}$ is valid when positive.
If every diagonal is nonpositive, successive endpoint replacement does
not increase the objective. Endpoint dynamic programming is then exact,
with at most two labels per coordinate, without a growth assumption.

The quantitative analysis assumes a unique optimizer $x^*$ and
\begin{equation}\label{eq:growth}
F(x)-\OPT\ge g\norm{x-x^*}^2\quad(x\in X),\qquad
\OPT=F(x^*),\qquad \kappa=\max\{1,L/g\}.
\end{equation}
This is global growth in the original Euclidean coordinates. Near-tied
integer assignments may make $g$ small even if the continuous extension is
strongly convex. Rescaling coordinates can change $L/g$ and is not a
free invariance of the theorem.

The lower-bound arguments first use exact factor evaluations and exact
arithmetic. The bit-complexity theorem applies to explicit rational
quadratics; Section~\ref{sec:polynomial} states its polynomial extension.
An arbitrary real-valued oracle does not supply a bit-complexity theorem.

\section{A corrected grid is a continuous certificate}\label{sec:bound}

Work in a current product subbox $B\subseteq X$. For each coordinate take
a finite feasible grid $G_i$ containing both current endpoints. At a node
$v\in G_i$, let $\ell_i(v)$ be the largest adjacent interval length.
For integer coordinates ignore intervals of length one; set $\ell_i(v)=0$
if no other adjacent interval remains. Define
\begin{equation}\label{eq:correction}
d_i(v)=\frac L8\ell_i(v)^2,\qquad
D(y)=\sum_i d_i(y_i),\qquad Q(y)=F(y)-D(y),\qquad
b=\min_{y\in\prod_iG_i}Q(y).
\end{equation}
Assign each unary correction to one containing bag. This changes no
factor scope or graph edge.

\begin{lemma}[Interpolation and conditional interpolation]\label{lem:round}
For every $x\in B$, there is a distribution on grid assignments $Y$ with
$\E Y=x$ and $\E Q(Y)\le F(x)$. Thus $b\le\min_B F$.
Writing
\[
m_i(v)=\min\{Q(y):y_i=v,\ y\in\textstyle\prod_kG_k\},
\]
every adjacent pair $a,b'\in G_i$ satisfies
\begin{equation}\label{eq:conditional}
\min\{m_i(a),m_i(b')\}\le F(x)
\quad\hbox{whenever }x\in B,\ x_i\in[a,b'].
\end{equation}
If $y$ minimizes $Q$, then $F(y)-b=D(y)$.
\end{lemma}
\begin{proof}
Independently round each coordinate of $x$ to the endpoints of its
containing grid interval with probabilities preserving its mean. A grid
coordinate remains fixed. If an interval has endpoints $u,v$, its
rounding variance is $(x_i-u)(v-x_i)\le(v-u)^2/4$.
Concavity in~\eqref{eq:semiconcavity} gives, conditionally on all other
coordinates,
\[
\E F(x\hbox{ with coordinate }i\hbox{ rounded})
\le F(x)+\frac L2\operatorname{Var}(Y_i).
\]
Successive independent rounds keep these variances unchanged, so their
errors add. For a genuinely rounded interval of length $\Delta$, both
endpoint corrections are at least $L\Delta^2/8$, which dominates
$L\operatorname{Var}(Y_i)/2$. Integer unit intervals have no feasible
interior point, so omitting them causes no missing rounding error.
Coordinates already on their grids have zero variance and nonnegative
correction. Subtracting the corrections proves $\E Q(Y)\le F(x)$.
The minimum is no larger than this expectation. If $x_i\in[a,b']$,
the same rounding has $Y_i\in\{a,b'\}$, so its expected corrected value
is at least the left side of~\eqref{eq:conditional}. The last equality
is the definition of $b$ at its minimizing assignment.
\end{proof}

\subsection{One tree computation supplies all coordinate bounds}

Let $q_t$ be the sum of factors and negative unary corrections assigned
to bag $B_t$. For an oriented tree edge $t\to u$, let
$S_{tu}=B_t\cap B_u$. The message recurrence is
\begin{equation}\label{eq:messages}
M_{t\to u}(x_{S_{tu}})=
\min_{x_{B_t\setminus S_{tu}}}
\left[q_t(x_{B_t})+
\sum_{v\in\mathcal N(t)\setminus\{u\}}
M_{v\to t}(x_{B_v\cap B_t})\right].
\end{equation}
All variables range over their coordinate grids. Empty separators have
one scalar state. Two tree passes compute both orientations. Set
\begin{equation}\label{eq:calibrated}
A_t(x_{B_t})=q_t(x_{B_t})+
\sum_{v\in\mathcal N(t)}M_{v\to t}(x_{B_v\cap B_t}).
\end{equation}
Induction on a cut tree edge proves that $M_{t\to u}$ minimizes exactly
the factors in the $t$ component for the fixed separator labels. Running
intersection makes the remaining component variables disjoint once those
labels are fixed. Hence $A_t$ is the exact global corrected min-marginal
at bag $t$. Its minimum is $b$, and minimizing it with $x_i=v$ gives
$m_i(v)$ for any containing bag. Stored argmin choices recover a global
minimizer.

If $K=\max_i|G_i|$, all messages and coordinate min-marginals require
\begin{equation}\label{eq:dpwork}
O\bigl(p(N+|\mathcal A|)K^p\bigr)
\end{equation}
table operations, plus factor evaluation and index costs. Incoming
messages are accumulated by sums, and outgoing sums use exclusions;
there is no product over the children of a bag. This is the familiar
finite-state junction-tree computation, used here for an explicit
continuous lower bound. Variable elimination, tree agreement and
reparameterization are established tools; see, for example,
\cite{Dechter1999,WainwrightJaakkolaWillsky2005}.

\subsection{Filtering and the certificate for the original domain}

Maintain a feasible point realizing a nonincreasing upper bound $U$.
After finding $y\in\arg\min Q$, replace $U$ by $\min\{U,F(y)\}$.
Retain a coordinate interval $[a,b']$ exactly when
\begin{equation}\label{eq:retain}
\min\{m_i(a),m_i(b')\}\le U.
\end{equation}
Replace each coordinate domain by the hull of its retained intervals;
a singleton stays a singleton. Use $y$ as the next center.

\begin{proposition}[Safe filtering history]\label{prop:filter}
Every global optimizer, the incumbent, and $y$ survive the filter.
The last grid bound, together with all preceding filtering records,
is a valid lower bound on the original box. This assertion uses neither
growth nor uniqueness.
\end{proposition}
\begin{proof}
The conditional bound~\eqref{eq:conditional} makes every interval
containing a coordinate of a global optimizer or the incumbent pass
\eqref{eq:retain}. Also $m_i(y_i)\le Q(y)=b\le\OPT\le U$, so every
coordinate of $y$ survives. Taking hulls can only reintroduce points.
Any point removed from the original box has a first stage at which one
of its coordinates is excluded. At that stage its containing discarded
interval has conditional lower bound greater than the stage's $U$.
Its objective is therefore greater than the final nonincreasing $U$.
The last grid lower bound applies to all surviving points, and is at
most $\OPT\le U$. Together the two statements bound every original
point. Keeping only the final restricted grid would omit the first part
of this argument and is insufficient.
\end{proof}

The certificate records input factors, curvature evidence, grids,
corrections, directed messages or enough data to recompute them,
coordinate filtering decisions, and feasible upper bounds. A checker
can recompute the finite minima exactly. Growth estimates influence
neither the checking rule nor its validity.

\section{Why filtering removes the accuracy exponent}\label{sec:rate}

At stage $j$ use $h_j=s2^{-j}$ and a center $c$ in the current box.
For a fixed $0<\theta\le1/4$, build outward from $c_i$. At distance $t$
take a continuous step $h_j+\theta t$ or an integer step
$\max\{1,\lfloor h_j+\theta t\rfloor\}$, clipping at an endpoint.
The relevant adjacent length satisfies
\begin{equation}\label{eq:mesh}
\ell_i(v)\le h_j+\theta|v-c_i|.
\end{equation}
For an integer interval of length at least two this follows from the
floor step; unit intervals have zero correction. Clipping only decreases
the length. Thus for every grid assignment $z$,
\begin{equation}\label{eq:penaltyenergy}
D(z)\le\frac L4nh_j^2+
\frac{L\theta^2}{2}\left(\norm{z-x^*}^2+\norm{c-x^*}^2\right).
\end{equation}

\begin{lemma}[Contraction]\label{lem:contraction}
Suppose $\theta^2\le1/(8\kappa)$. Put
$B_0=\max\{1,4L/(11g)\}\le\kappa$. If $y_j$ is the stage's corrected
minimizer, and the initial center is any feasible point, then
\begin{align}
E_j:=\norm{y_j-x^*}^2
&\le\frac{4L}{15g}nh_j^2+\frac{E_{j-1}}{15},\label{eq:contract}\\
E_j&\le B_0nh_j^2,\qquad
F(y_j)-b_j=D(y_j)\le\frac{7L}{8}nh_j^2.\label{eq:stagegap}
\end{align}
Here $E_{-1}$ denotes the initial center's squared error.
\end{lemma}
\begin{proof}
The restricted box still contains $x^*$. Lemma~\ref{lem:round} and
growth give $gE_j\le F(y_j)-\OPT\le D(y_j)$. Since
$L\theta^2/2\le g/16$, absorbing the first distance term in
\eqref{eq:penaltyenergy} proves~\eqref{eq:contract}.
Initially $E_{-1}\le ns^2$, and
$4L/(15g)+1/15\le B_0$. At later stages $h_{j-1}=2h_j$, so induction
uses $4L/(15g)+4B_0/15\le B_0$. At all stages
$E_{j-1}\le4B_0nh_j^2$. Substitution into~\eqref{eq:penaltyenergy}
gives
\[
D(y_j)\le\frac L4(1+10\theta^2 B_0)nh_j^2
\le\frac{7L}{8}nh_j^2,
\]
where $\theta^2B_0\le1/4$ is a deliberately loose bound. Filtering
does not alter any step of this proof, because it retains the optimizer
and the next center.
\end{proof}

\begin{lemma}[Localization of retained intervals]\label{lem:localize}
After stage $j$, the retained continuous domain in coordinate $i$ lies
within distance $5\sqrt{n\kappa}\,h_j$ of $(y_j)_i$. The integer domain
lies within distance $1+5\sqrt{n\kappa}\,h_j$.
\end{lemma}
\begin{proof}
An interval passing~\eqref{eq:retain} has an endpoint $v$ and a complete
grid assignment $z$ with $z_i=v$ and $Q(z)=m_i(v)\le U_j$.
By~\eqref{eq:stagegap}, $U_j-\OPT\le7Lnh_j^2/8$. Consequently
\[
g\norm{z-x^*}^2
\le\frac{9L}{8}nh_j^2+
\frac{L\theta^2}{2}\bigl(\norm{z-x^*}^2+4B_0nh_j^2\bigr).
\]
Absorb as before to obtain
\begin{equation}\label{eq:witnessradius}
\norm{z-x^*}^2\le
\left(\frac{6L}{5g}+\frac{4B_0}{15}\right)nh_j^2
\le\frac{22}{15}\kappa nh_j^2.
\end{equation}
The qualifying endpoint is within
$(\sqrt{22/15}+1)\sqrt{n\kappa}\,h_j$ of $(y_j)_i$.
Its continuous interval length is at most
$h_j+\theta(\sqrt{22/15}+2)\sqrt{n\kappa}\,h_j$, by comparison
with the previous center. Since $\theta\sqrt\kappa\le1/\sqrt8$
and $n,\kappa\ge1$, their sum is below
$5\sqrt{n\kappa}\,h_j$. Taking the hull preserves this radius.
An integer interval with positive correction obeys the same bound.
A unit interval may survive through its good endpoint while its other
endpoint is bad; its physical length contributes the stated extra one.
\end{proof}

\begin{lemma}[Uniform state bound]\label{lem:states}
Under the hypotheses of Lemma~\ref{lem:contraction}, every stage has at
most
\begin{equation}\label{eq:statecap}
K(\theta,n)=100\theta^{-1}\lceil\log_2(n+2)\rceil
\end{equation}
nodes per coordinate, independent of its stage number and accuracy.
\end{lemma}
\begin{proof}
Stage zero has at most three nodes per coordinate since $h_0=s$.
For continuous grids, the successive unclipped distances are
$t_k=h((1+\theta)^k-1)/\theta$. A side of radius $R$ therefore needs
at most $1+\lceil\log(1+\theta R/h)/\log(1+\theta)\rceil$ steps.
For integer grids put $H=\max\{h,1\}$. Every nonfinal integer step is
at least $(H+\theta t)/3$: if $h\ge1$, use
$\lfloor u\rfloor\ge u/2$; if $h<1$ split at $\theta t=2$.
Thus $t+H/\theta$ grows by a factor at least $1+\theta/3$.
The same count holds with $H$ and $\log(1+\theta/3)$.

At the next stage $h=h_j/2$. Lemma~\ref{lem:localize} bounds
$\theta R/h$ in the continuous case and $\theta R/H$ in the integer
case by quantities below $1+8\sqrt n$. Since
$\log(1+\theta/3)\ge\theta/4$ for $\theta\le1/4$, summing both
side counts and the center is bounded by~\eqref{eq:statecap}.
The constant 100 includes ceilings and is not an optimized practical
recommendation.
\end{proof}

\subsection{A capped schedule needs no supplied growth constant}

For $\varepsilon>0$, choose the smallest nonnegative integer $J$ with
\begin{equation}\label{eq:stagebudget}
\frac{7Lns^2}{8}\,4^{-J}\le\varepsilon.
\end{equation}
This uses rational comparisons, not logarithm oracles. Try
$\theta=2^{-\mu}$ for $\mu=2,3,\ldots$, restarting from the original
box and its lower endpoint vector as center each time. Run stages zero
through $J$. While generating a grid,
abort the trial before any coordinate exceeds~\eqref{eq:statecap},
before allocating that trial's bag tables. Stop only when the verified
global gap $U-b$ is at most $\varepsilon$.

The first trial with $\theta^2\le1/(8\kappa)$ has
$2^\mu\le6\sqrt\kappa$, passes the cap, and succeeds by
Lemma~\ref{lem:contraction}. Earlier trials remain sound even when their
guessed resolution is unsuitable. The caps bound their work by a geometric
sum of state spaces. Without caps an earlier unsuccessful trial could
pay an accuracy-dependent power before the successful one was reached.

The decisive absorption is
\begin{equation}\label{eq:logabsorb}
\lceil\log_2(n+2)\rceil^p\le(C_0p)^p(n+2)
\end{equation}
for an absolute $C_0$. It follows from
$\sup_{t\ge0}t^pe^{-t}=(p/e)^p$, adjusting the logarithm base and
ceiling. A power of the accuracy bits could not be absorbed this way;
the state bound has removed precisely that dependence.

\begin{theorem}[Certified quadratic approximation]\label{thm:approx}
For an explicit rational mixed-integer box quadratic with supplied
decomposition and~\eqref{eq:growth}, the capped algorithm returns a
rational feasible point and an independently checkable gap at most
$2^{-q}$ in
\begin{equation}\label{eq:fpt}
f(p,\kappa)(I+q+1)^C
\end{equation}
bit operations, where $C$ is absolute. Neither $g$ nor $\kappa$ is an
input. Certificate validity needs neither growth nor uniqueness.
\end{theorem}
\begin{proof}
The preceding lemmas prove correctness and bound state counts and stages.
It remains to bound arithmetic. After fixed-coordinate substitution,
choose a common positive denominator $D$ for coefficients, endpoints,
and $L$. Its bit length is polynomial in $I$.
For a fixed trial let $K$ be its cap, allowing one extra candidate for
abort detection. At stage $j$, the unclipped continuous displacement is
\[
h_j\frac{(2^\mu+1)^k-2^{\mu k}}{2^{\mu(k-1)}}.
\]
All grid coordinates have denominators dividing
$A_j=D2^{j+\mu K}$. Indeed inherited centers and endpoints have the
previous common denominator, which divides this one; new displacements
have the displayed denominator; clipping selects an inherited endpoint;
integer floors stay integral. Denominators do not multiply over stages.

All quadratic values, corrections, and message entries have denominators
dividing $8DA_j^2$. Their numerators have bit length polynomial in
$I+j+\mu K$, since coordinates stay in the original box and each message
is a sum of a subset of factor and correction values at some assignment.
Factor evaluations, message arithmetic, comparisons, backtracking,
certificate storage and checking have the same polynomial bit bound.
The stage budget is $O(I+q+1)$. Combine~\eqref{eq:dpwork},
\eqref{eq:statecap},~\eqref{eq:logabsorb}, and
$2^\mu=O(\sqrt\kappa)$, summing the earlier capped trials. All factors
depending on $p,\kappa$ belong to $f$, and the remaining polynomial
exponent is absolute.
\end{proof}

Approximation also terminates without growth. For a sufficiently fine
trial $\theta\le h_J/s$, each corrected interval at stage $J$ is at
most $2h_J$, so the gap is at most $Lnh_J^2/2\le4\varepsilon/7$.
Before that stage continuous grids have at most $s/h_j+3$ nodes, and
integer grids at most $2s/h_j+3$ when $h_j\ge1$ or $s+2$ otherwise.
These fit the cap for sufficiently small $\theta$. This observation
gives termination, but not the growth-based rate.

\section{Exact output is a separate theorem}\label{sec:exact}

\subsection{Rational quadratics}

Write the quadratic, after substituting fixed coordinates, as
\[
F(x)=(x^{\mathsf T}Hx+h^{\mathsf T}x+k)/D,
\]
where $H$ is symmetric integral, $h,k$ are integral, and all endpoints
have denominators dividing the positive integer $D$. Include factors of
two from off-diagonal normalization when selecting $D$. Put
\begin{equation}\label{eq:height}
C_H=\max\{1,\max_{ij}|H_{ij}|\},\quad
\Delta=(2nC_H)^n,\quad R=D\Delta,\quad V=DR^2.
\end{equation}
These numbers have polynomial bit length; they are not enumerated.

\begin{lemma}[Rational height]\label{lem:height}
The problem has a global optimizer with a common coordinate denominator
at most $R$, and its optimum value has reduced denominator at most $V$.
Every isolated global optimizer satisfies the coordinate bound.
\end{lemma}
\begin{proof}
Fix the integer assignment of a global optimizer. Among optimizers in
that slice select one whose smallest containing continuous-box face has
minimum dimension. Write $S$ for its free continuous coordinates and
$A$ for integer and bound-active coordinates; $x_A=t_A/D$ with integral
$t_A$. The free Hessian is positive semidefinite by local optimality.
If singular, a free null direction has both zero first-order change and
zero quadratic change. Move along it until a free coordinate hits a box
bound. The objective stays optimal and the face dimension decreases,
a contradiction. Thus $2H_{SS}$ is positive definite, and stationarity is
\[
(2H_{SS})(Dx_S)=-Dh_S-2H_{SA}t_A.
\]
Cramer's rule gives common denominator $D\det(2H_{SS})$, bounded by
$D|S|!(2C_H)^{|S|}\le R$; an empty determinant is one.
Evaluating the integer-coefficient quadratic at a common-denominator
point gives reduced value denominator at most $DR^2=V$.
For an isolated optimizer a singular free Hessian would itself yield
a short segment of optimizers, so the same argument applies directly.
\end{proof}

\begin{lemma}[Uniqueness supplies qualitative quadratic growth]
A quadratic with a unique optimizer on a compact mixed box has some
$g>0$ satisfying~\eqref{eq:growth}.
\end{lemma}
\begin{proof}
Otherwise a sequence $x_j\ne x^*$ has ratio
$(F(x_j)-F(x^*))/\norm{x_j-x^*}^2\to0$. Compactness and uniqueness
force $x_j\to x^*$, and hence their integer coordinates eventually
equal those of $x^*$. Write $t_j=\norm{x_j-x^*}$ and take a subsequence
with $u_j=(x_j-x^*)/t_j\to u$, $\norm u=1$. The exact quadratic
expansion is
\[
\frac{F(x_j)-F(x^*)}{t_j^2}
=\frac{\nabla F(x^*)^{\mathsf T}u_j}{t_j}
+\frac12u_j^{\mathsf T}\nabla^2F\,u_j.
\]
First-order optimality in the fixed integer slice makes the first term
nonnegative. Its boundedness implies $\nabla F(x^*)^{\mathsf T}u=0$,
and the limiting quadratic term is nonpositive. The limit direction is
in the continuous-box tangent cone with zero integer components, so
$x^*+\tau u$ is feasible for all sufficiently small positive $\tau$.
Optimality forces the same quadratic term nonnegative. It is zero,
and the entire short ray is optimal, contradicting uniqueness.
\end{proof}

This qualitative lemma does not bound conditioning usefully. A small
input can encode an exponentially weak growth margin.

\begin{theorem}[Exact rational quadratic output]\label{thm:exact}
Under the hypotheses of Theorem~\ref{thm:approx}, the unique optimizer
and exact optimum value can be found and certified in
$f_1(p,\kappa)(I+1)^{C_1}$ bit operations for absolute $C_1$, without
receiving $g$. Every unique rational mixed-box QP has such a positive
$g$, although its conditioning may be poor.
\end{theorem}
\begin{proof}
Distinct rationals with denominators at most $K$ differ by at least
$1/K^2$. An interval $[b,U]$ of width at most $1/(4V^2)$ therefore
isolates the unique possible optimum value of height at most $V$.
Continued fractions reconstruct this rational in polynomial bit time.
Around each coordinate of the feasible incumbent use a reconstruction
window of radius $1/(4R^2)$. Its length is below $1/R^2$, so there is
at most one rational of denominator at most $R$ inside it. If all
coordinates reconstruct, check the original bounds, integrality, and
exact equality of the candidate's objective to the isolated value.
Accept only after those checks. Acceptance is globally valid even before
the incumbent is known to be close.

Call the approximation algorithm with $\varepsilon=2^{-q}$ for
$q=1,2,4,8,\ldots$. When
\[
\varepsilon\le\min\{1/(4V^2),\ g/(32R^4)\},
\]
growth places the incumbent within $1/(\sqrt{32}R^2)<1/(4R^2)$ of
$x^*$. Every coordinate reconstructs correctly and the check succeeds.
The required $q$ is $\poly(I)+O(\log\kappa)$, using $g\ge L/\kappa$;
doubling overshoots by at most two. The approximate bit bound and
polynomial reconstruction cost prove the claim. The height constants
come from the original problem, not from refined hull endpoints.
\end{proof}

Only the optimal value is guaranteed to have denominator at most $V$.
Arbitrary continuous feasible values do not occupy a fixed value lattice.
The reconstruction and equality checks cannot be replaced by rounding
an arbitrary small objective gap to zero.

\begin{proposition}[Exact acceptance using the candidate's denominator]
\label{prop:candidateheight}
Suppose $b\le\OPT$ is a certified original-domain lower bound and $x$
is feasible. Write $F(x)=a/W$ in lowest terms, with $W>0$. If
\[
F(x)-b<\frac1{VW},
\]
then $x$ is globally optimal. This implication requires neither uniqueness
nor a growth constant.
\end{proposition}
\begin{proof}
The optimum is rational with reduced denominator at most $V$ by
Lemma~\ref{lem:height}. If $F(x)>\OPT$, the difference of the two
rationals is at least $1/(VW)$. Since $b\le\OPT$, that contradicts
the displayed strict inequality.
\end{proof}

This rule uses the denominator of the actual candidate; it does not
assume that a feasible value has denominator at most $V$. It allows the
implementation to verify candidates from reconstruction or exact
stationarity recovery independently of how they were generated. A large
candidate denominator can demand a very small certified gap, and the
rule alone supplies no running-time bound for finding an accepted point.

\begin{theorem}[Finite exact recovery without uniqueness]
\label{thm:generalfinite}
Every nonempty bounded rational mixed-box QP admits finite exact output
by the default geometric conditioning schedule, endpoint snapping and
rational stationary-face recovery. The implementation succeeds with
sufficiently large finite round, stage, table, LP-pivot and time limits.
No growth constant or optimizer is supplied. This statement gives no
general polynomial or fixed-parameter running-time bound.
\end{theorem}
\begin{proof}
The height bound $R$ can also bound vertices of every stationary
polytope formed by fixing original integer labels and bounds and imposing
a subset of free-gradient equations. In scaled coordinates $u=Dx$, a
vertex matrix consists of integer stationary rows and unit bound rows.
Expansion along the latter leaves a possibly nonprincipal Hessian minor.
The same determinant estimate used in~\eqref{eq:height} bounds such a
minor too. Thus each vertex has a common coordinate denominator at most
$R$, including its original-bound slacks. The implementation's sharper
row-product bound has the same nonprincipal-minor property.

If no coordinates remain, direct evaluation is exact. Otherwise $n\ge1$.
Let $S$ be the entire optimal set, $\tau=1/(4nR)$, and suppose a feasible
point $y$ is within $\tau/2$ of a nearest optimizer $s$. Their integer
labels agree. Snap each continuous coordinate within $\tau$ of an
original endpoint to that endpoint. Original endpoint separations are
at least $1/D>2\tau$, so the rule is unambiguous and snaps every
coordinate active at $s$ correctly. Let $J_0$ be the interior continuous
coordinates at $s$, and $J\subseteq J_0$ those left free after snapping.

Let $P$ retain the original box, fix the labels and active coordinates
of $s$, and impose $\nabla_{J_0}F=0$. This nonempty bounded polytope
has constant objective $F(s)$: differences of its points are supported
on $J_0$, and subtracting their gradient equations annihilates the free
Hessian action. Let $q(x)$ sum the nonnegative slacks of the additional
endpoints selected in $J_0\setminus J$. Then
\[
q(s)\le3n\tau/2=3/(8R)<1/R.
\]
A positive minimum of $q$ at a vertex of $P$ would be at least $1/R$,
by the common denominator bound. Therefore some $s'\in P$ meets all
selected endpoints simultaneously. The implemented selected-face
stationarity system is feasible. Every feasible solution $x$ of that
system has the same value as $s'$: for $d=x-s'$, both free gradients
vanish, $d$ is supported on $J$, and $H_{JJ}d_J=0$. Exact quadratic
expansion gives $F(x)=F(s')=\OPT$, without a nonsingularity assumption.

The unconditional approximation argument after
Theorem~\ref{thm:approx} reaches every positive accuracy with sufficiently
large finite caps. Compactness implies that feasible objective values
tending to $\OPT$ approach $S$ in distance. Thus sufficiently accurate
incumbents satisfy the preceding recovery premise. Exact LP finds a
recovery point even for singular faces, and a gap below $1/V^2$ makes
Proposition~\ref{prop:candidateheight} accept it. Each successful finite
prefix uses finite resources. The proof supplies existence of such caps,
not a useful bound on them for arbitrary optimal sets.
\end{proof}

This specializes the earlier stationary-polytope recovery mechanism to
the implemented denominator bound. It returns one rational optimizer
and value. A compact representation of every optimizer, or an efficient
width-and-set-growth bound for discovering an unknown optimal set, is
a separate question.

\subsection{Explicit polynomial factors}\label{sec:polynomial}

Suppose the factors are explicit rational monomial lists of fixed degree
$d\ge2$. A valid rational curvature certificate can be computed by
differentiating twice in each coordinate, evaluating the resulting
monomials by interval arithmetic on the full continuous hull, and summing
their upper endpoints. A sharper checked bound may replace this one.
Its certificate must have polynomial verification cost; include the
accepted $L$ and any such certificate in $I$. The quantitative parameter
uses the bound actually accepted.

\begin{corollary}[Polynomial approximation and native-integer exactness]
At fixed numerical degree $d$, under~\eqref{eq:growth}, the pruned-grid
algorithm returns a certified $2^{-q}$ approximation in
$f_d(p,\kappa)(I+q+1)^C$ bit operations with exponent independent of
$p,\kappa$. If every coordinate is a native integer, it also returns an
exact optimizer in $f_d(p,\kappa)(I+1)^C$ bit operations.
\end{corollary}
\begin{proof}
The interpolation, filtering, and growth proofs use coordinate
semiconcavity, not a quadratic expansion. For arithmetic, polynomial
values at stage $j$ have denominators dividing $DA_j^d$, and values
including corrections have a common denominator $8DA_j^{\max\{d,2\}}$.
Explicit fixed-degree evaluation is polynomial in input and coordinate
bit length, which proves the approximation statement.

For native integers choose a common denominator $Q_0$ of the original
objective coefficients. Every feasible value lies in $Q_0^{-1}\Z$.
A certified gap less than $1/Q_0$ at a feasible incumbent excludes every
strictly better integer value. It suffices to take
$q=1+\lceil\log_2Q_0\rceil=O(I)$. The checker verifies feasibility,
the original objective value, coefficient denominators, and the complete
pruning certificate. Growth is unnecessary for this stopping rule's
validity.
\end{proof}

The result does not automatically give exact continuous polynomial
coordinates. For example $x^3-6x$ on $[1,2]$ has optimizer $\sqrt2$ and
value $-4\sqrt2$. Nor does uniqueness imply quadratic growth for
polynomials: $x^4$ at zero is a counterexample. Arbitrarily large
binary-encoded exponents or general circuits also need a separate
evaluation-size analysis. A rational lattice rescaling can alter the
conditioning parameter, so the native-integer theorem is not silently a
metric-preserving lattice theorem.

\section{Extensions and the remaining structural questions}\label{sec:extensions}
\subsection{Negative curvature: a useful metric does not yet give sparse states}
For a continuous-box quadratic with Hessian $H$, let
$\nu=\max\{0,-\lambda_{\min}(H)\}$. A stronger desired parameter would
replace $L/g$ by $\max\{1,\nu/g\}$, allowing large positive curvature
to affect encoding length only. The target is continuous: $\nu=0$ in
a mixed-integer problem still permits convex integer optimization.

Two existing identities explain both the promise and the missing step.
For $d=x-x^*$, first-order optimality on the convex box gives
$F(x)-\OPT\ge d^{\mathsf T}Hd/2$. Combining it with growth, with
weights $g/(g+\nu)$ and $\nu/(g+\nu)$, yields
\begin{equation}\label{eq:convexenergy}
F(x)-\OPT\ge\frac{g}{2(g+\nu)}
d^{\mathsf T}(H+2\nu I)d.
\end{equation}
Thus the exact positive energy has the desired conditioning. A dense
change of coordinates to make this energy Euclidean generally destroys
the box factorization. Independent rounding need not preserve its
cancellations. The inequality itself does not bound separator states.

Likewise, for $\lambda>\nu$ define the convex-recourse envelope
\[
E_\lambda(c)=\min_{x\in X}
\{F(x)+\lambda\norm{x-c}^2/2\}.
\]
Subtracting $\lambda\norm c^2/2$ makes it a pointwise infimum of affine
functions of $c$, hence concave. Its upper coordinate curvature is at
most $\lambda$. Minimizing the lower bound
$g\norm{x-x^*}^2+\lambda\norm{x-c}^2/2$ over all $x\in\R^n$ gives
\[
E_\lambda(c)-\OPT\ge\frac{g\lambda}{2g+\lambda}\norm{c-x^*}^2.
\]
The ratio becomes $2+\lambda/g$. Each rational evaluation is a strongly
convex quadratic program, but the envelope can be dense even when $H$
is tridiagonal. Efficient evaluation is not bounded-width factorization.

\paragraph{A conditional-moment obstruction.}
Keeping exact positive energy is useful but not sufficient by itself.
For lifted moments with mean $m\in X$ and covariance
$\Sigma=Y-mm^{\mathsf T}\succeq0$, the relaxed objective satisfies
\[
\tfrac12\langle H,Y\rangle+b^{\mathsf T}m+c
=F(m)+\tfrac12\langle H,\Sigma\rangle
\ge F(m)-\tfrac\nu2\operatorname{tr}\Sigma.
\]
If these moments are conditioned on a single global cell of widths $w_i$,
interval moment bounds give $\Sigma_{ii}\le w_i^2/4$ and loss at most
$\nu\sum_iw_i^2/8$. A mixture indexed by complete global cell choices
inherits the same weighted bound. Matching finitely many separator
moments in individual bag cells does not ensure this representation.

An explicit witness uses $(s,u_1,u_2,v)\in[0,1]^4$, rational
$0<h\le1/2$, and
\begin{align*}
F_h={}&8(s/h-(u_1+u_2)/2)^2+8(s/h-1/4-v/2)^2\\
&+u_1(1-u_1)+u_2(1-u_2)+v(1-v)+(u_1+u_2+v)/16.
\end{align*}
The bags are $\{s,u_1,u_2\}$ and $\{s,v\}$. With
$\delta=s/h-1/8-(u_1+u_2+v)/4$, expansion gives
\[
F_h-\tfrac14=16\delta^2+
2[(1-v)u_1u_2+v(1-u_1)(1-u_2)]+(u_1+u_2+v)/16.
\]
This proves unique optimum $(h/8,0,0,0)$, value $1/4$, and growth
$g=1/22$: the squared distance is at most
$2\delta^2+11\norm{(u_1,u_2,v)}^2/8$. The Hessian is PSD squares
plus diagonal negative curvature $-2$ on controls; direction
$(0,1,-1,0)$ attains eigenvalue $-2$, so $\nu/g=44$.

For $t=s/h$, put left-bag masses $1/8,3/8,3/8,1/8$, respectively, at
\[
(t,u_1,u_2)=(0,0,0),\ (1/2,1,0),\ (1/2,0,1),\ (1,1,1).
\]
Put equal right-bag masses at $(t,v)=(1/4,0),(3/4,1)$.
Both separator laws have moments $\E t=1/2$, $\E t^2=5/16$,
$\E t^3=7/32$. All local squares and concave penalties vanish, giving
relaxed value $3/32$ and gap at least $5/32$. On mesh $2h$, every
separator atom lies in its first cell and controls occupy endpoint cells;
all occupied widths tend to zero. Even a global PSD completion exists:
with common mean $(h/2,1/2,1/2,1/2)$ take
\[
\Sigma=aa^{\mathsf T}/16+3bb^{\mathsf T}/16,
\quad a=(h,1,1,2)^{\mathsf T},\quad b=(0,1,-1,0)^{\mathsf T}.
\]
Its bag restrictions match the measures. Thus a width-squared error bound
depending only on $p,\nu/g$ fails for this particular local-moment
interface. It is not a lower bound against adaptive filtering or all
moment hierarchies; matching the fourth separator moment already excludes
this witness. The companion supplies the complete construction and
checks, including global McCormick inequalities.

\subsection{A complete reduction for certified affine convex recourse}
There is a useful class on which stiff positive curvature can be removed
exactly. Retain $z$ in a rational product box and let private, mutually
disjoint blocks $y_t\in[\ell_t,u_t]$ attach only to coordinates $E_tz$.
For one block, suppress $t$ and write
\[
f(y,z)=\tfrac12y^{\mathsf T}Cy+y^{\mathsf T}Dz+c^{\mathsf T}y
+\tfrac12z^{\mathsf T}Az+d^{\mathsf T}z+e,\qquad C\succeq0.
\]
Supply affine maps $\bar y(z)=a+Bz$, $\lambda(z)$, and $\mu(z)$
that satisfy, throughout the attachment box,
\begin{align*}
&\ell\le\bar y(z)\le u,\qquad\lambda(z),\mu(z)\ge0,\\
&C\bar y(z)+Dz+c=\lambda(z)-\mu(z),\\
&\lambda_i(z)(\bar y_i(z)-\ell_i)=0,
\qquad\mu_i(z)(u_i-\bar y_i(z))=0.
\end{align*}
Affine box ranges, coefficient identities, and rational PSD elimination
check these conditions in polynomial bit time, including singular $C$.
Put $q_t(z)=f_t(\bar y_t(z),z)$. Quadratic expansion and
complementarity give the exact nonnegative identity
\begin{equation}\label{eq:affinerecourse}
f(y,z)-q_t(z)=\tfrac12(y-\bar y)^{\mathsf T}C(y-\bar y)
+\lambda(z)^{\mathsf T}(y-\ell)+\mu(z)^{\mathsf T}(u-y)\ge0.
\end{equation}
Equality holds at $\bar y(z)$. Thus minimizing the original objective
is exactly minimizing $q(z)=q_0(z)+\sum_tq_t(E_tz)$, whose block
Hessian is $A+B^{\mathsf T}CB+B^{\mathsf T}D+D^{\mathsf T}B$.

If the original problem has point growth $g$, evaluating that inequality
on the affine response yields
\begin{equation}\label{eq:affinemetric}
q(z)-\OPT\ge g(z-z^*)^{\mathsf T}
\left(I+\sum_tE_t^{\mathsf T}B_t^{\mathsf T}B_tE_t\right)(z-z^*).
\end{equation}
In particular Euclidean growth $g$ transfers. If every attachment fits
in a supplied residual bag of size at most $p$, reduced factors retain
that width. If all reduced Hessian diagonals are nonpositive, endpoint
DP is exact without growth. Otherwise set
$L_{\rm red}=\max_i(\nabla^2q)_{ii}>0$; Theorem~\ref{thm:exact}
applies with this reduced upper coordinate curvature and
$\kappa_{\rm red}=\max\{1,L_{\rm red}/g\}$; the affine map lifts its
exact output. Identity~\eqref{eq:affinerecourse} transfers every lower
bound and certificate to the original problem independently of growth.
When $L_{\rm red}\le C_0\beta$, for a checked rational
$\nu\le\beta<2\nu$ and fixed $C_0$, this is a scoped algorithm
parameterized by $p$ and $\max\{1,\nu/g\}$. The condition on reduced
positive curvature remains an additional hypothesis.

Fixed-active-set affine recourse is standard multiparametric quadratic
programming \cite{BemporadEtAl2002}. The contribution here is its checked
composition with the conditioned sparse algorithm and metric growth
transfer. The relevant width is the residual width after elimination;
original width alone does not control fill-in. The companion note also
gives a ladder with arbitrarily many negative eigenvalues, fixed reduced
conditioning, and arbitrarily stiff penalties $M\sum_i(y_i-u_i)^2$.
Its algebraic certificate demonstrates the parameter separation without
claiming a difficult benchmark. General changing active faces and
their curvature cancellations require further structure; a general
negative-curvature-only bound is still open.

The global affine map need not be supplied for a fixed private-block
partition. Choose the strict interior midpoint $z_0$ of the attachment
box and solve one exact central convex QP, obtaining $\widehat y$.
All central minimizers have the same gradient
$h=C\widehat y+Dz_0+c$: PSD expansion between two minimizers forces
their difference into $\ker C$. Coordinates with $h_i>0$ must be at
their lower bounds, those with $h_i<0$ at upper bounds. Any globally
affine feasible selector fixes these coordinates at those bounds,
because an affine bounded function attaining a bound at an interior
parameter is constant. For $h_i=0$, the selector's gradient must vanish
identically. If its coordinate is not identically a bound, it is interior
at interior parameters, and stationarity applies; if it is constant at a bound, its
affine gradient has one sign and vanishes at the midpoint, so is zero
everywhere. Thus existence is equivalent to a polynomial-size linear
feasibility system for affine selector coefficients: these forced bounds,
zero gradient identities, full-box primal bounds and the remaining
gradient signs. Absolute-value auxiliary variables encode affine box
ranges. One exact central QP and one LP therefore decide existence and
construct the rational certificate, including singular PSD blocks. A
failed test rejects this affine-selector class, not the underlying
optimization problem. The companion gives the full recognition proof;
choosing a useful private-block partition remains outside it.

\paragraph{Changing active sets without explicit value-function pieces.}
A second exact composition keeps private convex blocks as value factors.
Let $Y_t$ be fixed nonempty bounded rational polytopes, independent of
the retained variables, and define
\[
\psi_t(v)=\min_{y\in Y_t}
\{\tfrac12y^{\mathsf T}C_ty+y^{\mathsf T}D_tv+c_t^{\mathsf T}y\},
\qquad C_t\succeq0.
\]
For every fixed $y$ the parameter dependence is affine, so $\psi_t$ is
concave. In $V(z)=q_0(z)+\sum_t\psi_t(E_tz)$, the upper coordinate
curvature is bounded by the direct quadratic diagonals of $q_0$.
Every rational value query is an exact convex QP with primal and rational
KKT certificates \cite{KozlovTarasovKhachiyan1980}. The attachment scopes
fit a supplied residual
decomposition; the number of retained variables need not be small.
Original full-vector growth implies retained-coordinate growth by
evaluating it at conditional minimizers.

The corrected-grid proof now applies with local convex-QP evaluations
instead of polynomial factor evaluation. Different active sets can give
different denominators, but each message is a selected sum of at most
the input number of polynomial-bit oracle values. Its denominator bit
length is bounded by their summed lengths, so the same parameterized
approximation bound holds. Exact recovery uses height bounds for the
\emph{original} rational QP on a bounded polytope, not for the nonsmooth
value function. Reconstruct the retained vector, lift through the exact
oracles, and accept only at the isolated original optimum value.
The companion proof establishes the full bit and height argument.
This handles changing active faces but need not cancel large direct
positive curvature: for $M(y-z)^2$ the direct diagonal is $2M$.
Affine condensation can first remove blocks with a global response
certificate, leaving other blocks as value factors.

\paragraph{Local pieces can preserve curvature cancellation.}
The affine and changing-active-set cases admit a stronger common interface.
For a private convex block write
\[
f(y,v)=\tfrac12y^{\mathsf T}Cy+y^{\mathsf T}Dv+c^{\mathsf T}y
       +\tfrac12v^{\mathsf T}Av+d^{\mathsf T}v+e,
\qquad C\succeq0.
\]
Suppose a rational binary space partition covers its attachment box.
Every closed leaf carries an affine feasible response $\bar y=a+Bv$
and affine nonnegative bound multipliers satisfying stationarity and
complementarity throughout that leaf. A partition tree proves coverage;
rational LP dual witnesses prove the affine inequalities. The identity
\[
f(y,v)-f(\bar y(v),v)
=\tfrac12(y-\bar y)^{\mathsf T}C(y-\bar y)
 +\lambda(v)^{\mathsf T}(y-\ell)+\mu(v)^{\mathsf T}(u-y)\ge0
\]
then verifies the reduced quadratic on that entire leaf. Distinct response
maps may agree only in value at a shared boundary, which is sufficient.
Let $H_r=A+B_r^{\mathsf T}CB_r+B_r^{\mathsf T}D+D^{\mathsf T}B_r$
be the reduced Hessian on leaf $r$ and put $\ell_j=\max_r(H_r)_{jj}$.

\begin{proposition}[Curvature across a certified partition]
\label{prop:piecewisecurvature}
The conditional value $\phi(v)=\min_y f(y,v)$ has upper coordinate
curvature $\ell_j$ in coordinate $j$. Therefore local partitions for
independent blocks give the valid reduced coordinate bound
\[
L_i=(H_0)_{ii}+\sum_{t:i\in E_t}\max_r(H_{tr})_{j(t,i),j(t,i)},
\]
where $H_0$ is the retained quadratic's Hessian. No global overlay of
the local partitions is needed.
\end{proposition}
\begin{proof}
Restrict a coordinate line. The finite partition makes $\phi$ continuous
and piecewise quadratic, with second derivative at most $\ell_j$ within
each piece. Moreover $\phi(v)-v^{\mathsf T}Av/2-d^{\mathsf T}v-e$
is an infimum of affine functions of $v$, hence concave. Its derivative
jumps on the line are downward. Adding or subtracting a smooth quadratic
does not change those jumps. Consequently the derivative of
$\phi-\ell_jv_j^2/2$ is nonincreasing within and between pieces.
The same argument holds on a partition boundary, where overlapping
pieces agree as restricted polynomials. Summing the coordinate bounds
proves the displayed formula.
\end{proof}

Let $S$ include the encoding lengths of all partition and KKT witnesses.
With unique full-vector growth $g$ and residual maximum bag size $p$,
the corrected-grid algorithm now has the previous bounds
$f(p,\max\{1,L/g\})\poly(S+q)$ for approximation and
$f_1(p,\max\{1,L/g\})\poly(S)$ for exact rational output, where
$L=\max_i L_i>0$. Indeed, a factor query traverses one local tree,
original growth implies retained growth, and exact output uses the
original QP height. If every $L_i\le0$, endpoint DP is exact without
growth. Certificate size is part of the input to this statement;
constructing a small partition from an arbitrary QP is not guaranteed.

For a concrete separation, take $M\ge1$, $0\le z\le3/4$ and
$0\le y_1,y_2\le1$, with
\[
f_M=M(y_1+y_2-z)^2+(y_1-2z+1/2)^2.
\]
The three parameter intervals have reduced values
$(1/2-2z)^2$ on $[0,1/4]$, zero on $[1/4,1/2]$, and
$M(z-1/2)^2/(M+1)$ on $[1/2,3/4]$. Their curvature maximum is $8$,
whereas the direct retained diagonal is $2M+8$. The unique conditional
response changes slope, so no globally affine response represents it.
The companion note extends this to a ladder with arbitrarily many
negative eigenvalues, uniformly bounded negative curvature and full-vector
growth, and three pieces per block. It supplies that family's complete
parameter proof and the finite construction by active patterns. The
implemented constructor and independent checker cover scalar attachments;
the general certified-partition theorem does not assert an implemented
automatic multidimensional partition algorithm. Parametric critical
regions themselves are classical \cite{BemporadEtAl2002}.

\subsection{Why a global grid correction cannot simply be made local}
For $m=32$ and $\epsilon=1/16$, on $[0,1]^{m+1}$ consider
\[
F(x,y)=x^2-\frac x2\sum_{i=1}^m y_i+\sum_{i=1}^m y_i^2
+\left(\frac m{16}-1-\epsilon\right)x.
\]
This is a star with bags $\{x,y_i\}$ and coordinate curvature $L=2$.
Exact leaf recourse has $y_i=x/4$ and value
$V(x)=-x^2+(1-\epsilon)x$, so its unique optimizer has
$x=1$, $y_i=1/4$, and value $-1/16$.
On the grid $\{0,1/2,1\}$ the conditional grid values at the three
center labels are $0$, $23/32$, and $31/16$. The grid incumbent is zero.
For the bag cell $x\in[1/2,1]$, $y_i\in[0,1/2]$, which contains the
optimizer's bag projection, the smallest corner grid min-marginal is
$23/32$. A proposed bag correction $|B|Lh^2/8=1/8$ would reject the
cell. The missing outside error is
\[
V_h(x)-V(x)=m\dist(x/4,\{0,1/2,1\})^2,
\]
which is not constant in the separator and does not disappear under
message normalization. This refutes that local substitution, not every
possible conditional-recourse method.

\subsection{The precise conditional-recourse interface}
For a bag $B$, define the true original-domain value
$W_B(v)=\min_{x_{-B}}F(v,x_{-B})$. When the outside feasible set is
independent of $v$, pointwise infima preserve coordinate semiconcavity:
after subtracting $Lv_i^2/2$, each member is concave, and the infimum
of concave functions is concave. For a bag cell of side at most $h$,
rounding only its bag coordinates gives lower bound
\[
\min_{v\text{ a corner}}W_B(v)-e_B,
\qquad e_B=|B|Lh^2/8.
\]
No outside variables are rounded.

Suppose an oracle gives $\ell_B(v)\le W_B(v)\le u_B(v)$, a feasible
completion attaining $u_B(v)$, and $u_B(v)-\ell_B(v)\le\eta\le ae$,
where $e=pLh^2/8$ is a common budget. First update the incumbent from
all current-level completions and then retain cells whose corner lower
bound minus $e$ is at most $U$. Interpolation at an optimizer gives
$U\le\OPT+e+\eta$. Every retained cell has a corner with
\[
W_B(v)\le\OPT+2(e+\eta).
\]
This is the sufficient local-error interface. Obtaining these bounds
requires certified optimization over the original outside domain;
ordinary grid messages do not satisfy the interface automatically.

\subsection{An exact cut oracle for a large nonconvex residual}
The conditional interface has a concrete realization. For
\[
F(x)=c+\sum_i(a_ix_i^2+b_ix_i)+\sum_{i<j}c_{ij}x_ix_j,
\]
supply a core $C$ of $k$ continuous variables and residual $R$ such that
$a_i\le0$ for $i\in R$ and signs $s_i\in\{-1,1\}$ satisfy
$c_{ij}s_is_j\le0$ for residual edges. A graph traversal checks the
sign consistency for a supplied core. There are no restrictions on edges
touching the core, and the residual graph may be dense.

Fix a rational core assignment. Successive coordinatewise concave
endpoint choices show that a residual optimum exists at endpoints.
For integer residual coordinates round their bounds inward first; the
same endpoint argument applies without enumerating their integer labels.
Choose signed endpoint representations $x_i=\alpha_i+d_iy_i$ with
$y_i\in\{0,1\}$ and $d_i=s_i(u_i-\ell_i)$. The residual energy is
\[
K+\sum_i h_iy_i+\sum_{i<j}q_{ij}y_iy_j,
\qquad q_{ij}=c_{ij}d_id_j\le0.
\]
With $r_{ij}=-q_{ij}/2$ and $p_i=h_i+\frac12\sum_{j\ne i}q_{ij}$,
this equals $K+\sum_ip_iy_i+\sum_{i<j}r_{ij}|y_i-y_j|$.
Represent $y_i=1$ on the source side of a cut. Arcs $i\to t$ of
capacity $\max\{p_i,0\}$, $s\to i$ of capacity
$\max\{-p_i,0\}$, and opposite pair arcs of capacity $r_{ij}$ give
\begin{equation}\label{eq:cutidentity}
F=K+\sum_i\min\{0,p_i\}+\operatorname{cut}(y).
\end{equation}
A feasible flow of the same value as the returned cut proves conditional
optimality by weak duality. Rational max-flow is polynomial in query
encoding length, and the identity, flow and cut need only linear-size
checks after graph construction. The class is stable under residual
subboxes and added linear terms.

The cut reduction is classical \cite{KolmogorovZabih2004}; its role here
is a certified outside oracle that removes the entire residual dimension
from discretization. Separate concavity is essential to this route.
General continuous quadratic submodularity is not being treated as an
exact cut oracle. In particular the current semidefinite-representation
paper proves exactness in dimensions at most three and exhibits a
four-dimensional gap \cite{BurerNatarajanWillemsen2026v3}.

The companion mixed-submodular result broadens this residual class when,
after sign reversal, the nonpositive-diagonal variables can be sent to
endpoints and the remaining continuous principal block is PSD. Minimizing
that convex block gives a submodular set function of the endpoint labels:
apply the lattice inequality to minimizing completions and their
coordinatewise minimum and maximum. Exact submodular minimization and
convex-QP queries then give polynomial-time recourse. Its polynomial
certificate uses a convex combination of at most $|R|+1$ greedy base
vectors, certified convex-QP chain values, and matching primal value.
The greedy-base certificate is also classical
\cite{IwataFleischerFujishige2001}; the added contract is exact
convex-QP witnesses and box-stable composition.
The full constructive certificate proof remains in the companion note.
The completed reference backend implements rational cutting planes for the
Lov\'asz extension, conditional convex-QP queries and the resulting
certificate. Its finite capped implementation is distinct from the
polynomial-time submodular oracle used in the complexity theorem.

For core domain $[0,1]^k$, query corner values of
$V(v)=\min_{x_R}F(v,x_R)$ on nested dyadic core cells. With core upper
coordinate curvature $L$, a side-$h_j$ cell has error
$e_j=kLh_j^2/8$. Update the incumbent using all queried completions,
retain cells with $\min_{\rm corners}V-e_j\le U$, and generate only
children of retained cells. This has a complete trace certificate and
gap at most $e_j$, by the interface proved above.

There is also a precise smoothed count. Perturb core linear coefficients
independently and uniformly on $M$ equally spaced points of
$[-\sigma,\sigma]$, with $M\ge\max\{2,2^J\}$ chosen before sampling.
For a fixed level-$j$ tuple with value at most $\OPT+2e_j$, comparison
with both coordinate neighbors confines an interior coordinate's noise
to an interval of length at most $Lh_j+4e_j/h_j=Lh_j(1+k/2)$.
The interval depends on the unperturbed conditional value and the fixed
tuple, not on other core noises. Its probability is at most its length
divided by $2\sigma$, plus $1/M$. Summing the independent product over
all full-grid tuples bounds their expected number by
\[
H_k=\left[3+\frac{(1+k/2)L}{2\sigma}\right]^k.
\]
Every retained cell touches such a tuple; a tuple touches at most $2^k$
cells, each has $2^k$ children and each child $2^k$ corners. Hence the
expected number of conditional queries through level $J$ is at most
$2^k+8^kJH_k$. Multiplying by polynomial query and certificate cost
gives an absolute-exponent polynomial in residual size and precision,
with the exponential dependence confined to $k$ and $L/\sigma$.
Correctness holds on every draw, and the optimized objective is the
perturbed objective. This realizes the conditional interface on a
supplied-core class; it does not solve the corresponding question on
arbitrary sparse graphs.

\paragraph{Implemented exact quadratic output.}
This oracle also permits deterministic exact output with no growth
assumption, although the cell count can be large. Let $S_0$ be the product
of denominators of objective coefficients and original box endpoints,
$T_0=S_0^3$, and $H_0\ge1$ the largest absolute entry of the integer
scaled core Hessian $T_0\nabla^2_{CC}F$. For any residual endpoint label,
$T_0F$ is an integer-coefficient quadratic in the unit-box core. A
smallest-face optimizer has nonsingular free Hessian, as in
Lemma~\ref{lem:height}. Its determinant is at most $k!H_0^k$, so its
value denominator is at most
\[
D_0=T_0(k!H_0^k)^2.
\]
The same bound applies to the global value without enumerating all
residual labels. Run core search to gap $1/(2D_0^2)$. Fix the incumbent's
residual label and solve its core QP by enumerating its $3^k$ faces and
nonsingular stationary systems, including vertices. The minimum-face
argument makes this finite list complete. Its value $\widehat v$ lies
between $\OPT$ and $U$, and both $\widehat v$ and $\OPT$ have
denominators at most $D_0$. Rational separation forces equality.
The checker only verifies the search trace, height bound, candidate
feasibility, candidate value denominator at most $D_0$, and that value's
membership in the narrow interval; it need not repeat face
enumeration. The implemented mode reports an unfinished enclosure if its
level budget is insufficient. This stopping depth must not simply be
inserted into the finite-noise count: sampled coefficients affect $D_0$,
and requiring $M\ge2^{J(D_0(M))}$ could be circular. The separate exact
smoothed theorem uses an existing base-selected closure and fallback
construction, not that substitution.

\subsection{Coupled constraints and nonunique optima}
The product-domain hypothesis does mathematical work. Independent
rounding may violate an equality even when the equality is sparse.
Constraint extensions must supply a feasible rounding law, an exact
reduction, or recourse with a verified error budget. A repair map must
also account for distortion of objective error, conditioning, and
integrality. Merely placing constraint scopes in small bags does not
complete this argument.

For a nonunique optimal set $S$, set growth
$F(x)-\OPT\ge g\dist(x,S)^2$ need not localize all retained coordinates
near one center. The point-growth proof above therefore cannot be
relabelled as a set-growth proof. Special structure or an explicit
selection mechanism is needed. Exact polynomial boundary output likewise
requires a certificate appropriate to active faces and algebraic
coordinates; quadratic rational reconstruction is not such a certificate.

\subsection{A complete extension for integral TU fibers}
Let
\[
X=\{(x,z):\ell\le x\le u,\ z_i\in Z_i,\ Ax+Bz\le b\},
\]
where $A$ is totally unimodular (TU), $\ell,u,B,b$ are integral, and
each $Z_i$ is an explicitly listed finite integer set. Constraint rows
and objective factors fit supplied bags of maximum size $p$. Let $n_c$
count nonfixed continuous coordinates, $d=\max_i|Z_i|$, and
$W=\max_i(u_i-\ell_i)$. Assume a verified full directional bound
$\nabla^2_{xx}F\preceq LI$ on the continuous hull, uniformly over
finite labels, and unique-point growth $g$ on $X$.

At level $j$ use aligned mesh $h=2^{-j}$, keeping endpoints of retained
continuous hulls and all retained finite labels. For fixed feasible
$(x,z)$ write its surrounding cell as $x=h(k+t)$. The fiber is
\[
At\le2^j(b-Bz)-Ak,\qquad0\le t\le1.
\]
Its right side is integral and adjoining box bounds preserves TU.
Therefore every vertex is binary, and $t$ is a convex combination of
feasible binary corners. This is classical box integrality
\cite{ChervetGrappeRobert2018}. The resulting feasible random corner
$Y$ has $\E Y=x$ and $\E\norm{Y-x}^2\le n_ch^2/4$. Taylor's upper
bound gives
\begin{equation}\label{eq:tuerror}
\E F(Y,z)\le F(x,z)+E_j,\qquad E_j=n_cLh^2/8.
\end{equation}
Correlations in this rounding explain why a diagonal-curvature bound
alone would not suffice.

Exact tree DP rejects infeasible bag assignments and computes the least
feasible grid objective $v_j$ and coordinate min-marginals $M_i$.
Then $b_j=v_j-E_j$ is a lower bound. The current incumbent lies on each
later nested grid, so after updating it $U=v_j$. Retain a continuous
interval $[s,t]$ when $\min\{M_i(s),M_i(t)\}-E_j\le U$, and a
finite label $t$ when $M_i(t)-E_j\le U$. Conditional feasible rounding
proves these exclusions sound and gives the same full-history
certificate argument as Proposition~\ref{prop:filter}. An empty initial
grid certifies infeasibility by the same rounding lemma.

A retained endpoint or label has a full feasible grid witness with
value at most $U+E_j\le\OPT+2E_j$. Growth puts it within
$r_j=\sqrt{n_cL/(4g)}h$ of the optimizer. Continuous retained hulls
therefore have length at most $2r_j+2h$, and the next grid has at most
$5+2\sqrt{n_cL/g}$ nodes per continuous coordinate. Thus table work is
\begin{equation}\label{eq:tucost}
O\!\left(p(N+|\mathcal A|+\#\mathrm{rows})
 [K_0^p+(J+1)K^p]\right),
\end{equation}
where $K_0=\max\{d,W+1\}$,
$K=\max\{d,5+\lceil2\sqrt{n_c\kappa}\rceil\}$, and the least
$J$ with $E_J\le2^{-q}$ is $O(I+q+1)$. Dyadic coordinates and
fixed-degree explicit evaluation have polynomial arithmetic overhead.
No growth constant is used by this algorithm or its checker.

For rational quadratics, finite-label fixation and shrinking continuous
hulls allow exact rational reconstruction. On a minimum-dimensional
optimal continuous face, the Hessian is positive definite on the face
tangent space by the null-direction argument of Lemma~\ref{lem:height}.
The integral KKT saddle matrix has dimension at most $2n_c$ and bounded
entries, so Cramer's rule gives a polynomial-bit common denominator
bound. Reconstruct only after hulls isolate bounded-height coordinates
and the value interval isolates a bounded-height value; then check all
original constraints and the exact objective. The companion TU note
supplies explicit determinants and thresholds. The level count becomes
$\poly(I)+O(\log\kappa)$ with the same state factors.

This is a complete accuracy-independent state theorem, but it is
\emph{not} FPT in $(p,\kappa)$: the factor $n_c^{p/2}$ remains.
Large integral capacities incur the initial pseudopolynomial $K_0^p$
cost. Network conservation with unit capacities and integral finite-label
switch contributions is a concrete application; every conservation row
still contributes its full scope to width. General nonlinear constraints,
unaligned rational fibers, and compressed large finite-label sets are
outside this statement.

Two reviewed refinements keep the original variables and bag scopes.
First, with original equalities $Cx+Ez=a$, the curvature bound is needed
only for directions in $\ker C$, because feasible rounded corners have
the same labels and preserve those equalities. For a quadratic, a rational
basis of $\ker C$ or its orthogonal projector verifies the restricted
bound without substituting variables into objective factors. A penalty
$\rho\norm{Cx+Ez-a}^2$ then changes neither grid values nor rounding
allowances. Coefficient-based exact height bounds can still grow with
$\rho$, so this statement concerns per-level search and correction.
Second, a supplied positive rational initial mesh $\delta$ is admissible
when $\ell/\delta,u/\delta,(b-Bz)/\delta$ are integral for all listed
labels. Check one reference label vector and each individual column-label
increment, without enumerating their Cartesian product. The initial
count becomes $W/\delta+1$, while the filtered count is unchanged and
$h_j=\delta2^{-j}$. Large common units can therefore remove the initial
capacity cost; arbitrary unaligned magnitudes still retain that cost.

\paragraph{Nonunique TU minima and unions of retained cells.}
Exact output can be extended to arbitrary nonunique TU quadratics, and
a stronger state representation helps when each continuous coordinate
has only finitely many optimal values. Let $S=\arg\min_XF$ and use the
quadratic error bound $F(v)-\OPT\ge g\dist(v,S)^2$ for some $g>0$.
Such a constant exists on each compact quadratic/polytope slice
\cite{LuoSturm2000}; finitely many other integer slices have positive
value gaps. No useful numerical lower bound on $g$ is asserted.

Keep the union of surviving mesh cells, including singleton nodes,
instead of their interval hull. Every feasible point still lies in one
mesh cell within the retained domain, so the same TU corner distribution
and error $E_j$ apply. In particular
\[
b_j=v_j-E_j\le\OPT\le U_j=v_j\le\OPT+E_j.
\]
A passing endpoint has a feasible witness of value at most
$\OPT+2E_j$, and hence lies within
$a_j=\sqrt{n_cL/(4g)}h_j$ of some optimal coordinate value. Every point
of a retained cell lies within $a_j+h_j$ of such a value. If the projection
$S_i$ has at most $r$ elements, the next mesh therefore has at most
\[
r\bigl(5+\lceil2\sqrt{n_c\kappa}\rceil\bigr)
\]
labels in coordinate $i$. This replaces the continuous part of $K$ in
\eqref{eq:tucost}. Different coordinates may have witnesses near different
optimizers; the argument neither enumerates nor chooses optimal
combinations. There can be exponentially many isolated optima. Positive-
dimensional continuous optimal sets are excluded by the finite-projection
state bound, but not by the exact-output statement that follows.

For exact output, the general-polytope recovery lemma in the companion
work computes original-data rational thresholds $\delta,\tau>0$ and
an optimum denominator bound $V$, all of polynomial binary length.
Given a feasible point within distance $\delta$ of $S$, fix its integer
labels, impose original inequality rows with slack at most $\tau$ as
equalities, and solve rational feasibility with stationarity along that
face. Every feasible recovery output is optimal. The lemma's proof uses
stationary-polytope heights and simultaneous slack snapping; this is an
imported recovery result, not a claim that arbitrary KKT points are global.
Its stationarity multipliers are unrestricted.

Refine until $E_j<1/(4V^2)$, isolate the bounded-height optimum value in
$[b_j,U_j]$, and try the recovery LP. Accept a point only after checking
all original constraints and equality to that isolated value. Early
recovery can fail or produce a nonglobal stationary point, but neither
can pass the final equality check. Eventually $E_j\le g\delta^2$, so
recovery succeeds, including for a continuum of optimizers. The required
level is $\poly(I)+O(\log\kappa)$; favorable grid-size bounds additionally
require finite projections. This extends exact termination without a
supplied growth constant. It does not resolve the stronger width-FPT
problem for arbitrary unknown optimal sets.

\subsection{Compact descriptions of nonunique optimal sets}
Two restricted results avoid pretending that all retained coordinates
converge to a single unknown optimum.

First suppose $F(x)=x^{\mathsf T}Qx+b^{\mathsf T}x+c$ on a mixed box
and $Q_{ii}\le0$ for every coordinate. Endpoint DP gives root value
$\OPT$. Root the decomposition and let $m_t$ denote its upward message.
For assigned bag factor $\phi_t$, define the nonnegative endpoint
Bellman residual
\[
r_t(v_{B_t})=\phi_t(v_{B_t})+\sum_{c\text{ child}}m_c(v_{S_c})
-m_t(v_{S_t})\ge0,
\]
where the root message is $\OPT$. On any endpoint assignment the messages
telescope to $F(v)-\OPT=\sum_tr_t(v_{B_t})$.
Let $W_t(v;x)$ be the product of endpoint interpolation weights
$(u_i-x_i)/(u_i-\ell_i)$ or $(x_i-\ell_i)/(u_i-\ell_i)$.
Independent mean-preserving endpoint rounding preserves cross terms and
adds variance $(x_i-\ell_i)(u_i-x_i)$ to $x_i^2$. Taking expectations
therefore gives the exact identity
\begin{equation}\label{eq:endpointset}
F(x)-\OPT=
\sum_t\sum_{v\in\{\ell,u\}^{B_t}}r_t(v)W_t(v;x)
-\sum_iQ_{ii}(x_i-\ell_i)(u_i-x_i).
\end{equation}
All terms are nonnegative. Optimality is equivalent to mixed feasibility,
vanishing of every positive-residual product, and endpoint status for
every $Q_{ii}<0$. This is a factored polynomial description of all
optimizers with at most $N2^p+n$ equations. It is constructed and checked
in $2^{O(p)}\poly(I)$ bit work with no growth promise and no enumeration
of interior integer labels or optimal components. Endpoint reduction
and Bellman reparameterization are classical; the displayed identity
records their full-set certificate.

Second, for a continuous-box quadratic
$F(x)=\tfrac12x^{\mathsf T}Hx+b^{\mathsf T}x+c$ and a feasible box
KKT point $s$, define
\[
d_i(s)=\begin{cases}
2\partial_iF(s)/(u_i-\ell_i),&s_i=\ell_i,\\
-2\partial_iF(s)/(u_i-\ell_i),&s_i=u_i,\\
0,&\ell_i<s_i<u_i,
\end{cases}\qquad D=\operatorname{diag}(d_i).
\]
These entries are nonnegative. Direct quadratic expansion yields
\begin{equation}\label{eq:diagonalcert}
F(x)-F(s)=\tfrac12(x-s)^{\mathsf T}(H+D)(x-s)
+\tfrac12\sum_id_i(x_i-\ell_i)(u_i-x_i).
\end{equation}
If $H+D\succeq0$, this certifies global optimality. The condition is
the classical box Lagrangian sufficient condition of
Li, Wu, and Quan \cite{LiWuQuan2015}, not a new certificate class.
It gives the full optimal set by
\[
\ell\le x\le u,\qquad(H+D)(x-s)=0,\qquad
d_i(x_i-\ell_i)(u_i-x_i)=0\quad\hbox{for all }i.
\]
If the test passes at one optimum, it passes at every optimum with the
same $D$: equality in~\eqref{eq:diagonalcert} puts $t-s$ in the
kernel and forces each positive-$d_i$ coordinate to an endpoint;
$\nabla F(t)=\nabla F(s)-D(t-s)$ then gives exactly the same entries.

The continuation combines this acceptance test with the earlier proximal
grid candidate generator and exact selected-face recovery lemma. Doubling
untrusted guesses $K$ for $L/g$, running each generator for its prescribed
finite budget, and accepting only a rational KKT/PSD certificate is sound
under every guess. On instances with at least one certificate of this
form, the first $K\ge\max\{1,L/g\}$ produces a globally optimal
candidate, and the invariance just proved ensures acceptance whichever
optimal component was recovered. The resulting exact discovery and
full-set description take $f(p,\max\{1,L/g\})\poly(I)$ bit work under
set growth toward $S$. Neither $S$ nor $g$ is supplied. The candidate
generator and recovery proofs remain in the identified companion notes;
the acceptance and full-set identity are proved here. Outside this
certifiable class the procedure may remain inconclusive.

\subsection{Polynomial boundary output with an explicit margin}
There is a restricted way to remove the full-Hessian positivity premise
from exact implicit polynomial output. Assume unique-point growth and
strict complementarity at every active continuous original bound.
Let $\gamma>0$ be the smallest inward derivative there, and let a checked
rational $M\ge1$ bound all absolute continuous-Hessian row sums on the
original hull. Put
$B_\gamma=\lceil\log_2\max\{2,M/\gamma\}\rceil$ for analysis only.

After retained integer labels have become singletons, work within their
fixed slice. In a globally retained continuous box of maximum half-width
$r$ and midpoint
$c$, the bound $\partial_iF(c)-Mr>0$ certifies fixing a lower-bound
coordinate; the corresponding negative upper bound certifies an upper
face. Each substitution preserves the minimum on the retained box.
An optional exact Bernstein coefficient bound also permits weak signs,
preserving at least one optimizer. This is standard monotonicity
propagation \cite{ArayaTrombettoniNeveu2010}, now combined with sparse
containment certificates.

At $r\le\gamma/(4M)$ all active coordinates are found. The remaining
free Hessian at the optimizer is at least $2gI$, since two-sided free
directions and the growth inequality give that second-order bound.
A rational third-derivative bound $T$ lets an exact midpoint test
$H_B(c)-(Tr+\tau)I\succ0$ certify a strongly convex patch on the
selected face. Dovetail capped grid trials with
$\theta_\mu=2^{-\mu}$, $K_\mu=4^\mu$, and
$\tau_\mu=L/(4K_\mu)$; do not tighten grading merely because more
precision is needed. The first $K_\mu\ge8\kappa$ reaches a certified
patch in $\poly(I)+O(\log\kappa+B_\gamma)$ stages.
Allocating at phase $q$ at most $2^{q-\mu}$ bit steps to each
$2\le\mu\le q$ gives finite overall work
$f_d(p,\kappa)\poly(I+B_\gamma)$.

The output is a fixed integer assignment, a selected continuous face,
and a rational strongly convex polynomial subproblem whose unique
minimizer is an original global optimizer. It is an exact implicit
description, not a rational-coordinate claim. Full search scheduling,
integer-label fixation, and patch-certificate details are proved in the
boundary companion. The margin term is material: existing examples have
exponentially large $B_\gamma$ at fixed degree, width and conditioning.
Weakly complementary boundary cases without another certificate remain
outside this termination theorem.


\section{Initial implementation and historical evidence}\label{sec:computation}
This section preserves the first-release measurements and their original
source snapshots. The later implementation and its separately saved runs
are described in Section~\ref{sec:completion}; the figures below are not
reruns of the updated solver.

The original pruned-grid diagnostic implemented two-pass tree messages,
all coordinate min-marginals, geometric grids, interval filtering, and
hull recentering. Its saved exact-rational record covers nine fixtures,
70 stages, 26,589 bag-table states, and 526 pruned intervals. Separate
finite enumeration checked min-marginals on 21,034 assignments. These
figures describe that diagnostic, not the continuation's benchmark suite.
It did not implement the unknown-growth caps, rational reconstruction,
or an independent stored-history verifier. Its successes establish
finite regression evidence, not competitive general solver performance.

The continuation distinguishes three empirical questions: whether the
produced certificate independently checks, whether filtering reduces
state counts against comparable unfiltered methods, and whether its
total time and memory are useful against other solvers. An ordinary
floating-point solver result is useful as a timing baseline but is not
automatically an exact global certificate. Controlled synthetic fixtures
and public-library instances must be reported separately, with explicit
limits and incomplete runs preserved.

\subsection{The initial sparse solver and independent replay}
The initial standard-library Python reference implementation handled supplied
general tree decompositions, branching trees, empty separators, continuous
and native-integer domains, fixed coordinates, and disconnected factor
graphs. It implements corrected geometric grids, all coordinate
min-marginals, domain filtering, and the capped unknown-growth
approximation schedule. Coordinate-specific diagonal corrections sharpen
the common-$L$ bound. Endpoint reduction, bounded exact coordinate
descent, and an optional rational PSD/KKT presolve supplement the grid
algorithm. A saved result contains original-model bounds and a complete
history, including when a resource limit interrupts refinement.

The separate verifier does not call grid generation, dynamic programming,
filtering, coordinate descent, or the convex presolve. It shares rational
model parsing and objective evaluation, independently reconstructs factor
tables and Bellman equalities, checks all min-marginals and exclusions,
and re-evaluates feasible incumbents. The certificate proves the embedded
rational model; a caller must also match that model to the intended input.
It is an independently implemented checker, not a formally verified one.
Replay can cost as much as solving because it recomputes finite-table
arithmetic. This first release returned certified approximations and certain
exact presolve or finite-grid answers. General continuous rational
reconstruction was a theorem at that stage; the completion described below
adds its implementation.

\subsection{Bounded comparison on original objectives}
The main suite has ten small rational instances and two unmodified binary
QPLIB instances. The small set includes random signed paths, bands and a
branching tree, mixed and integer paths, flat and tied optima, and two
positive-curvature scales. Random cases do not have planted optimizers;
deliberate structural cases are labelled diagnostics. An independent exact
integer-slice and active-face enumeration gives the small reference values.

All twelve cases use absolute tolerance $1/50$; a random path and both
curvature cases also use $1/1000$. The resulting 15 configurations are
run with four methods. Each has a two-second solver limit, with an
external 14-second worker limit including setup, serialization and replay.
Exact methods additionally have a 20,000 aggregate bag-state cap and
24-stage cap. SCIP is single-threaded with a 256 MB internal limit.
The grid methods share a greedy minimum-degree decomposition; SCIP does
not receive it. These timings come from one bounded run on a shared
machine, not a statistical performance study.

\begin{center}
\begin{tabular}{lrrl}
\toprule
Method & Small targets reached & Public targets reached & Incomplete runs\\
\midrule
Geometric, pruned & $13/13$ & $0/2$ & 2 table limits\\
Geometric, unpruned & $13/13$ & $0/2$ & 2 table limits\\
Uniform, unpruned & $11/13$ & $0/2$ & 4 table limits\\
SCIP (numerical bounds) & $13/13$ & $0/2$ & 2 time limits\\
\bottomrule
\end{tabular}
\end{center}

``Target reached'' means the requested objective tolerance, not exact
optimization. The SCIP small outcomes comprise six optimal statuses and
seven gap-limit statuses. Its dual bounds remain numerical; its repaired
box-feasible incumbents are separately evaluated exactly. All 45 rational
grid certificates, including incomplete runs, pass replay. All 39
small-instance grid enclosures contain their exact reference values.

For the 13 matched small configurations, the pruned geometric method
processes 6,167 cumulative bag-table states versus 27,386 without pruning.
The strongest reduction in this small set is the high-positive-curvature
case at tolerance $1/1000$: 771 versus 13,527 states. Tied-and-flat and
pure-integer cases have equal counts for the two geometric variants.
Thus pruning helps particular fixtures, while the suite does not show
universal pruning benefit. Median solve times on these 13 configurations
are approximately 0.0059 seconds with pruning and 0.0082 without it;
median replay times are 0.0084 and 0.0078 seconds, respectively. Median
whole-worker subprocess times are 0.174 and 0.193 seconds. Interpreter
startup, setup, exact arithmetic and checking materially exceed some solve
times. Peak worker memory, including replay, is about 25 MB on the small
grid runs; it is not isolated DP memory.

The public inputs are QPLIB\_3852 (231 binary variables, 440 quadratic
terms) and QPLIB\_5881 (120 variables, 2,123 terms). Their supplied
decomposition widths are 25 and 95, respectively, and are heuristic
upper bounds on treewidth. Each exact variant refuses its oversized
tables before allocation and retains a valid initial bound; SCIP reaches
its time limit. No public instance reaches the target. These outcomes
demonstrate a practical width barrier for the supplied representations;
they do not prove that better decompositions or formulations cannot work.
The benchmark validates QPLIB's triangular coefficient convention against
the bundled library solutions and preserves the original objectives and
domains, apart from sign reversal of a maximization objective.

\subsection{Dimension and certified affine reduction}
Fourteen additional runs test increasing path dimension, a width-two
case, and the affine-recourse reduction. Across both first-release suites,
all \textbf{56 rational certificates} pass independent replay; the four
affine enclosures also contain their analytic zero optimum. The main
and extension suites use 20.23 and 22.23 seconds of summed subprocess
wall time, respectively. No worker crashes or reaches its external wall
limit. The following timings keep solve and replay separate.

\begin{center}
\begin{tabular}{lrrrr}
\toprule
Pruned random path & Gap & Solve (s) & Replay (s) & Peak MiB\\
\midrule
$n=16$ & $0.01497$ & $0.0253$ & $0.0231$ & $24.3$\\
$n=64$ & $0.01361$ & $0.2020$ & $0.1967$ & $24.8$\\
$n=256$ & $0.05786$ & $2.0023$ & $1.2388$ & $46.5$\\
\bottomrule
\end{tabular}
\end{center}

The target is $1/50$. The first two paths finish; the 256-variable
pruned path hits its time limit. Its unpruned counterpart hits a table
limit at the same gap and uses 1.790 seconds for replay. SCIP reaches
a numerical gap of approximately $0.000612$ in 1.797 seconds on that
path. A pruned width-two, 16-variable case also reaches its target, with
gap $0.01953$, solve time $0.0318$ and replay time $0.0546$ seconds.
Three paths and one wider example do not constitute an empirical
complexity law, and width one does not ensure fast Python rational
arithmetic at every dimension.

The affine experiment has six original variables, four retained
variables, and an exactly checked identity
$F(u,v,y)=q(u,v)+M\sum_i(y_i-u_i)^2$ with feasible response $y=u$.
It is the companion's designed two-block example. The same reduced
rational core is solved once for all three stiffness values.

\begin{center}
\begin{tabular}{lrrrr}
\toprule
Representation & States & Solve (s) & Replay (s) & Gap\\
\midrule
Original, $M=1$ & $1{,}694$ & $0.0352$ & $0.0425$ & $0.000458$\\
Original, $M=100$ & $64{,}564$ & $2.0044$ & $2.0668$ & $0.399817$\\
Original, $M=1000$ & $76{,}986$ & $2.0033$ & $2.5630$ & $7.406978$\\
Reduced, all $M$ & $394$ & $0.0121$ & $0.0151$ & $0.000519$\\
\bottomrule
\end{tabular}
\end{center}

The target is $1/1000$. The stiff original cases time out; the reduced
case finishes and its rational incumbent and lower bound lift exactly
to each original instance. State counts are cumulative over completed
stages. This demonstrates a useful certified
preprocessing operation on the stated class. It does not show that
these analytically simple models are difficult for other solvers, and
the affine-selector recognizer itself is not benchmarked here.

\subsection{Initial extension implementations}
The first-release conditional-cut implementation included residual sign
recognition, an exact rational max-flow oracle, primal/flow certificates,
adaptive core cells, trace verification, and a rational-separation exact
mode. Its broader concave/convex submodular extension had small exact
diagnostics. The TU implementation was a bounded two-bag diagnostic;
affine recourse and optimal-set/boundary extensions also had exact
algebraic checks. General convex-QP factor evaluation, diagonal-class LP
recovery, and reusable constrained-tree search were implementation gaps
at that stage. Section~\ref{sec:completion} records which gaps the next
release closes and states the remaining limits.

The saved experiment files provide per-run times, peak process memory,
decomposition metadata, source hashes, rational reference values and
replay certificates. The report's coverage map points to them. The evidence
supports a reusable certifying reference solver and concrete extensions;
it does not establish competitive general-purpose MINLP performance.


\section{Completed implementation and new evidence}\label{sec:completion}
The second implementation stage addresses the first release's engineering
gaps. Its source and benchmark records are separate from the historical
snapshots in Section~\ref{sec:computation}. The mathematical distinction
between validity, termination and practical cost remains essential:
an accepted rational certificate is exact, whereas a resource-limited
search can finish without finding one.

\subsection{Shared finite tables and automatic decompositions}
The reusable finite-tree engine computes both message directions, all
coordinate min-marginals and attaining assignments for finite factor tables,
including infeasible entries. It checks tree structure and running
intersection; the model adapters check factor coverage. Its callers build
different objective and feasibility
tables but share this computational rule. Certificate checkers separately
validate the model-specific mathematical bounds. The base and constrained
checkers independently implement Bellman replay; the convex-value checker
reuses the finite-tree engine after independently checking its local
optimization witnesses.

When a box-QP decomposition is omitted, a deterministic minimum-fill
heuristic constructs one from the interaction graph. A supplied
decomposition remains supported and is checked. The resolved bags and tree
are serialized, so certificate replay does not depend on reproducing the
heuristic. Returned width is an upper bound, not a minimum-treewidth claim.
Table budgets are checked before oversized allocations. Sparse neighbor
lists reduce zero-product work in coordinate polishing and gradients;
dense model storage and exact rational arithmetic remain practical costs.

\subsection{Exact rational output in the general grid solver}
The \texttt{solve\_exact} wrapper repeatedly requests finer original-model
grid bounds, generates rational candidates, and applies
Proposition~\ref{prop:candidateheight}. The height bound is computed once
from the original coefficients and bounds. The implementation sharpens
\eqref{eq:height} using a product of integer Hessian row norms, which bounds
every relevant principal determinant by Hadamard's inequality. Neither
filtered hull endpoints nor candidate denominators alter that original
optimum-height bound.

Candidates include the current feasible incumbent, bounded-denominator
coordinate reconstruction, and exact stationarity on a proposed active
face. Singular stationary systems use rational linear feasibility.
These are candidate-generation procedures: stationarity or proximity
alone never establishes global optimality. The separate exact-output
checker replays the underlying original-domain certificate, verifies the
candidate against the original model, recomputes its objective and reduced
denominator, and checks the strict rational-separation inequality.

The wrapper exposes round, stage, table and time limits, with configurable
initial precision.
An interrupted run retains its certified enclosure and says that exact
output was not reached. Theorem~\ref{thm:generalfinite} now proves eventual exact output for
arbitrary nonunique rational box QPs with sufficiently expanded resource
caps. The stronger point-growth running-time theorem retains the
uniqueness hypothesis stated in Section~\ref{sec:exact}. General finite
termination does not provide that stronger efficiency guarantee.

\subsection{Reusable exact linear and convex quadratic optimization}
The shared rational optimization module implements two-phase simplex with
Bland's pivot rule. Successful LP results carry primal/dual witnesses,
infeasibility carries a Farkas witness, and unboundedness carries a feasible
point and improving ray. The result checker verifies the original rational
constraints and algebraic witness without replaying simplex pivots.
A pivot limit is inconclusive; it is not an infeasibility certificate.

The convex box-QP oracle checks positive semidefiniteness exactly and
searches active faces with exact linear algebra and LP feasibility for
singular faces. A feasible point with checked KKT signs certifies the global
minimum of the convex quadratic. Face and pivot limits bound this
reference implementation. Neither simplex nor active-face enumeration is
claimed here as a polynomial-time implementation of the theoretical
convex-QP oracle. Their purpose is to make the restricted algorithms
executable while retaining independently checkable answers.

\subsection{Original-model recourse composition}
The recourse pipeline recognizes globally affine convex responses,
substitutes accepted responses, solves the reduced model, and lifts the
result to the original coordinates. Each accepted response has a checked
positive-semidefinite private Hessian, feasible whole-box affine range,
and whole-box KKT identities and signs. Thus the original and reduced
optimum values coincide. The final certificate stores every reduction,
including fixed-coordinate substitutions, so a checker can reconstruct
the chain and verify the final incumbent in the original model.

Candidate private blocks can be supplied or proposed from continuous
components, private leaf-bag variables and single coordinates. This
selection is heuristic. A rejected block or exhausted recognition budget
remains in the model, and a failure to discover a reduction does not prove
that no useful partition exists. The reduced problem can use the sparse
grid solver or signed-concave minimum-cut recourse. The latter adapter
checks signs and separates a small core before applying the existing exact
conditional-cut search. The full reduction and search cost belongs in the
reported total; preprocessing is not a free input promise.

Changing-active-set convex recourse is also executable. Supplied disjoint
continuous PSD blocks, or heuristically discovered blocks, become local
conditional value factors. The implementation caches rational convex-QP
values and KKT witnesses at attachment grid points, uses corrected geometric
grids on the retained coordinates, and supports exact original-QP height
acceptance. Private feasible sets in this implementation are boxes; the
general private-polytope theorem has a wider input class. The checker
verifies local PSD/KKT data, recomputes finite-tree values using the shared
engine, and checks the original-domain history. It invokes no LP or QP
optimizer.

For scalar attachments, the integrated optional constructor enumerates
private active patterns within a stated budget, certifies a complete
interval partition, and uses its maximal private reduced curvature to
sharpen the retained diagonal. Contributions from multiple scalar blocks
add without a global overlay. The current constructor requires a positive-
definite private Hessian and nonfixed intervals; unsupported or over-budget
blocks retain the safe direct bound. The default pattern cap is $81$.
In tested changing-response examples this changes a retained bound from
$2M+6$ to $6$ for stiffness $M=1$ and $M=100$. General multidimensional
short-partition construction remains outside the implementation.

The mixed convex/concave submodular backend goes beyond minimum cuts.
After verified sign changes, endpoint reduction of separately concave
coordinates leaves a submodular set function evaluated by conditional
convex QPs. Rational LP cutting planes minimize its Lov\'asz extension;
greedy chains supply separating affine functions. The final lower bound
uses a finite convex mixture of checked greedy base vectors, and a feasible
conditional minimizer attains it. Its independent checker verifies the
conditional KKT witnesses and the mixture bound. This is an exact capped
reference implementation, without a polynomial iteration guarantee for
its cutting-plane search. It stores completed cached query witnesses, so
its certificate storage follows the query budget; the implementation does
not claim the theorem's polynomial certificate-size bound.


\subsection{Coupled constraints and optimal sets}
The reusable TU solver uses feasible common-mesh tables on general tree
decompositions, all coordinate min-marginals, and complete filtering
history. Constraint scopes join objective scopes in the decomposition.
The input includes an exactly checked TU certificate: supported structural
forms include signed incidence and consecutive-ones matrices; a bounded
all-minors fallback handles small matrices. A refused TU verification is
not a proof of non-TU structure. Original integral alignment, finite
integer labels and verified continuous curvature remain assumptions.
The optional union representation preserves separated surviving cells.
Exact output combines original-face rational candidates, the generic
stationary-face recovery LP and a height-separation check. With resource
caps removed, the recovery construction also covers arbitrary nonunique
quadratic minima; only the finite-projection subclass has the stronger
state bound. The width-dependent exponent and initial capacity costs
stated in Section~\ref{sec:extensions} are unchanged.

The optimal-set module constructs a compact Bellman-residual descriptor
for the entire optimal set of a coordinatewise concave mixed-box QP. A
separate checker verifies the recurrence and the support conditions,
including possible interior integer labels. For the diagonal-certifiable
continuous class, the implementation actually runs finite proximal trials
with guessed growth ratios, performs rational stationary recovery, and
accepts only the original KKT and positive-semidefinite certificate. A
false growth guess cannot justify acceptance. Exhausting the configured
search limits yields an inconclusive result, not a negative membership
decision for the class.

\subsection{Explicit polynomials and exact implicit boundary output}
The polynomial implementation accepts explicit rational monomial factors,
constructs or verifies their sparse decomposition, and computes upper
coordinate-curvature bounds by exact interval arithmetic. It uses the shared
finite-tree engine, corrected grids and complete pruning histories. Native
integer output has the coefficient-denominator lattice stopping rule from
Section~\ref{sec:polynomial}, computed after substituting any fixed rational
coordinates. Its checker independently recomputes that denominator and the
strict pre-acceptance gap. Continuous approximation does not round an
algebraic optimum to a rational exact answer.

The boundary search feeds finer globally certified polynomial boxes to a
separate exact acceptance test. It first fixes native labels, then
repeatedly certifies derivative signs and clamps the corresponding original
bounds. Such a clamp preserves at least one global optimizer; a weak sign
need not preserve every optimizer. On the remaining face, let $H_0$ be the
midpoint Hessian and let $\delta$ bound the maximum absolute row sum of
$H(x)-H_0$ throughout the box. An exact positive-definiteness check on
$H_0-\delta I$ proves uniform strong convexity. The polynomial together
with this rational face box is then an exact finite descriptor of one
global optimizer, possibly with irrational coordinates and value.

The boundary verifier independently recomputes monomial derivative
intervals, reductions and the matrix bound, after replaying the global
grid proof. It does not optimize the patch. The tested polynomial
$(x^2-1/2)^2+y+xy^2$ yields a certified patch for
$(x,y)=(1/\sqrt2,0)$ after four accuracy rounds. A bounded restart search
implements discovery; it does not implement the theoretical bit-step
dovetail or claim its FPT runtime. Failure of these sufficient tests, or a resource limit, can leave the
result inconclusive, and general growth-only polynomial boundary
output remains unresolved.

\subsection{New bounded comparisons and certificate checks}
The second-stage main suite has 73 configurations. Every solve runs
sequentially in a fresh process with one numerical-library thread, a
two-second cooperative solver budget, a five-second hard worker deadline,
and a 512 MiB address-space cap. Replay runs in a separate five-second
worker. The table cap is 30,000 states per stage. Saved timings separate
solving, replay, preprocessing and total subprocess time; the latter
includes both workers. Source snapshots and final hashes identify the
implementation actually measured. Earlier provisional runs are archived
and excluded from the following totals.

\begin{center}
\begin{tabular}{lrrr}
\toprule
Lane & Configurations & Requests completed & Valid proof replays\\
\midrule
Archived grid baseline & 16 & 12 & 16\\
Completed grid & 16 & 12 & 16\\
Exact rational output & 9 & 9 & 9\\
Integrated recourse & 11 & 10 & 11\\
TU constraints & 16 & 13 & 15\\
Optimal-set descriptions & 5 & 5 & 5\\
\midrule
Total & 73 & 61 & 72\\
\bottomrule
\end{tabular}
\end{center}

Completion means the request made in that lane: the stated approximate
gap, an exact value, infeasibility, or a checked optimal-set description.
Partial bounds can pass replay without completing the request. The one
configuration without a proof is an intentionally invalid TU input,
correctly rejected before solving. Across the main suite, 54 enclosures
contain independent exact reference values and 29 configurations certify
exact values. The nine exact-output runs cover eight distinct models;
a separable case repeated without convex presolve already has an exact
initial bound, so that repeat is not evidence of nontrivial reconstruction.
The fresh random path, band and mixed cases provide such evidence.

Both grid versions finish new random paths through 64 variables and a
32-variable random tree. They finish the same 12 of 16 requests, with
summed solve-plus-replay subprocess times 15.18 and 15.34 seconds. This
comparison jointly changes minimum-degree to minimum-fill decomposition,
the message engine and sparse polishing; it establishes neither a speedup
nor an isolated cause for a timing difference. Minimum fill reduces the
supplied widths of the two full binary QPLIB instances from 25 to 19 and
95 to 93. Both still exceed the table budget, and neither public input
reaches its target. The three continuous bound-only QPLIB entries screened
from metadata have 14,400--39,204 variables; one has unbounded coordinates.
They exceed the prototype's declared 512-variable screen. These metadata
refusals are not solver runs or successful continuous public benchmarks.

The integrated recourse tests operate on the original models. On a
shuffled 17-variable affine star, automatic preprocessing removes 16
coordinates and certifies value $-1$ in 0.065 seconds of solve time and
0.047 seconds of replay. The 33-variable dense minimum-cut example reaches
gap $1/1024$ in 0.089 seconds, with 0.099 seconds for replay. These are
controlled structural examples. The two changing-response examples with
stiffness $M=1$ and $M=100$ each use 17 table states, six levels and seven
conditional QP queries, with certified gap $27/65536$. Their identical
state counts illustrate this certified curvature reduction; two parameter
values do not establish an empirical complexity law.

Two same-model constrained comparisons expose the value of the added
structure. For $x_0(1-x_0)+(x_1-1/3)^2+(x_2-1/3)^2$ with $x_1=x_2$,
coordinate hulls hit the per-stage cap after 15 levels and 33,502 total
processed states. Union filtering returns exact value zero after 43
levels and 2,538 total states, using two stationary-recovery LPs and
24 pivots; exhaustive face search is disabled. For
$1000(x_0-x_1)^2+(x_0-1/3)^2$ with $x_0=x_1$, the ambient-curvature
mode times out after 252,192 total states. The certified equality-tangent
bound yields exact value zero with 1,938 states. These small examples
exercise specific representation and curvature choices, rather than
establishing general solver superiority.

Eleven separately frozen extension configurations exercise polynomial,
implicit-boundary and mixed-submodular functionality. Seven reach their
requested outcome; three retain checked intervals after a table limit,
a cut limit or an unsupported signed class. One bounded symmetric-quartic
boundary search remains inconclusive and emits no exact-output certificate.
All ten emitted certificates replay. Eight numeric intervals enclose
independent exact references, and two implicit patches contain independently
known analytic optimizers.

The non-growth polynomial $(x-1/3)^4$ reaches its requested gap in
36 stages and 545 states. A mixed polynomial needs eight stages and
469 states. The native-integer example, with one fixed rational coordinate,
certifies value $-53/6$ after two stages and 17 states: its raw gap $1/24$
is strictly below its value-lattice spacing $1/18$. Independent enumeration
checks all 100 native assignments. The unplanted six-variable signed mixed
QP returns $-509/105$ using nine conditional QP queries, five greedy cuts
and 88 pivots. A separate fixture requires a genuine mixture of two bases;
a single-base lower bound $-13/16$ improves to exact value zero.

Across both new suites there are therefore \textbf{84 configurations,
82 valid proof replays and 68 completed requests}. Fourteen other runs
retain checked incomplete or unsupported-class bounds. The invalid TU
input and inconclusive boundary search are the two outcomes without
certificates. Reference evidence comprises 62 numeric enclosures and two
implicit-optimizer containment checks. The extension workers use the same
hard time and memory limits and total 1.584 seconds of subprocess time;
no process reaches a hard deadline or address-space cap. These figures
exclude provisional archives and the unchanged 74 historical runs.

The final combined topic-specific test command passes 158 tests in
4.114 seconds. Its saved output and 35-source hash manifest accompany the
completion package. These tests cover solver contracts, replay, invalid
inputs, resource exits and certificate corruption. They are separate from
the bounded benchmark configurations, and neither is a project-wide or
CI verification claim.


\section{Evidence map and interpretation}\label{sec:evidence}
The corrected-grid proof, filtering-history argument, state bound,
unknown-growth schedule, rational arithmetic, quadratic height bound,
and exact-output construction appear in Sections~\ref{sec:bound}--
\ref{sec:exact}. They develop the existing October 2 pruned-grid result;
they are not claimed as newly discovered in this continuation.
Section~\ref{sec:extensions} includes the central structural and
composition proofs. Longer imported arguments are identified in the
companion coverage map. Sections~\ref{sec:computation}
and~\ref{sec:completion} distinguish historical measurements from the
completed implementation and its new runs.

\begin{longtable}{@{}>{\raggedright\arraybackslash}p{0.25\textwidth}
>{\raggedright\arraybackslash}p{0.36\textwidth}
>{\raggedright\arraybackslash}p{0.33\textwidth}@{}}
\toprule
Result & Established status & Boundary \\
\midrule
\endhead
Filtered mixed-box QP & Approximation and exact-output proofs;
implemented rational recovery and independent original-model replay &
Point-growth rate requires uniqueness; resource limits retain incomplete
certified enclosures \\
Rational optimization & Reusable exact LP and convex box-QP with
primal/dual, Farkas, ray or KKT witnesses & Simplex and active-face
implementations do not have the theoretical polynomial oracle's rate \\
Affine convex recourse & Exact reduction and growth transfer;
automatic candidate discovery, recognition and lifted certificates &
Discovery is heuristic; reduced fill and curvature must be controlled \\
Convex value factors & Exact composition and implemented conditional
box-QP factors with changing active sets & General private polytopes remain
theorem scope; direct retained curvature can remain large \\
Piecewise convex recourse & Reduced-piece curvature theorem without
a global overlay; integrated certified scalar construction & Explicit partition
size is part of theorem input; general short-partition discovery is not
promised \\
Conditional minimum cuts & Automatic residual/core recognition;
exact rational oracle, search, verifier and exact mode & Core size and
sign-consistent separate concavity; no generic width-only recourse \\
Mixed submodular recourse & Implemented greedy-cut LP search, conditional
QP witnesses and checked global mixture bound & Finite cut search and
active-face QP can be exponential; stored proof size follows query budget \\
Explicit polynomial grids & Rational interval curvature, sparse global
bounds and native-integer lattice exactness & Complexity theorem fixes
numerical degree; continuous exact output requires another certificate \\
Integral TU fibers & Feasible-grid theorem and reusable tree solver;
exact output with nonunique minima; union-state extension & XP width,
initial capacity/alignment costs; finite-projection state bound excludes
continuous optimal components \\
Optimal-set descriptions & Endpoint full-set constructor;
actual unknown-growth search and acceptance for diagonal class &
General efficient discovery for arbitrary unknown optimal sets remains
unresolved \\
Polynomial active face & Implemented exact implicit output with checked
face discovery & Termination theorem needs strict complementarity and an active-gradient
precision cost \\
General sparse $\nu/g$ target & Scoped reductions and
representation obstructions & No complete unrestricted algorithm \\
Practical solver value & Bounded comparisons and independent certificate
replays, with incomplete runs preserved & High-width refusals and exact
arithmetic costs prevent a broad competitive-performance claim \\
\bottomrule
\end{longtable}

A sound certificate, a quantitative running-time theorem, and measured
solver usefulness are different claims. Certificate validity follows from
checked rational inequalities and algebra. Runtime bounds require their
stated structural and growth assumptions and computational model.
Practical usefulness requires suitable experiments. Independent
research-agent review and finite checks do not constitute journal peer
review or a proof of priority.



\bibliographystyle{plain}
\bibliography{references}
\end{document}
